% ════════════════════════════════════════════════════════════════════════════
%  Reporte de Proyecto Final --- Compilador HULK
%  Asignatura: Compilación / Lenguajes de Programación
%  Facultad de Matemática y Computación, Universidad de La Habana
%
%  Estilo: inspirado en Classic Thesis (A. Miede) + Bringhurst, adaptado.
%  Compilar con: pdflatex REPORT.tex (dos pasadas para referencias)
% ════════════════════════════════════════════════════════════════════════════

\documentclass[11pt,a4paper,oneside,openany]{report}

% ── Codificación e idioma ───────────────────────────────────────────────────
\usepackage[utf8]{inputenc}
\usepackage[T1]{fontenc}
\usepackage[spanish,es-tabla,es-noindentfirst,es-noshorthands]{babel}
\usepackage{csquotes}

% ── Tipografía profesional ──────────────────────────────────────────────────
% mathpazo provee Palatino para texto y matemáticas: alta legibilidad,
% favorita en literatura clásica de informática (Knuth, Bringhurst).
\usepackage{mathpazo}
\linespread{1.08}
\usepackage{textcomp}
\usepackage{gensymb}

% ── Geometria de página ─────────────────────────────────────────────────────
\usepackage[
  top=2.8cm, bottom=2.8cm,
  inner=3.2cm, outer=2.4cm,
  headheight=15pt,
  headsep=18pt,
  footskip=24pt,
  marginparwidth=2cm
]{geometry}

% ── Color: paleta sobria y profesional ──────────────────────────────────────
\usepackage{xcolor}
\definecolor{halfgray}{gray}{0.55}
\definecolor{webgreen}{rgb}{0,.5,0}
\definecolor{webbrown}{rgb}{.6,0,0}
\definecolor{hulkblue}{HTML}{1F3A68}
\definecolor{hulkaccent}{HTML}{8B0000}
\definecolor{codebg}{HTML}{F8F8F8}
\definecolor{codeborder}{HTML}{CCCCCC}
\definecolor{codecomment}{HTML}{708090}
\definecolor{codekeyword}{HTML}{0033A0}
\definecolor{codestring}{HTML}{007F00}
\definecolor{codetype}{HTML}{6B3FA0}

% ── Encabezados de sección ──────────────────────────────────────────────────
\usepackage{titlesec}
\usepackage{titletoc}

% Capitulos al estilo Bringhurst: número grande gris, titulo en versalitas
% Estilo Bringhurst con itálica negrita (preserva tildes en Palatino)
\titleformat{\chapter}[display]
  {\bfseries\Large}
  {\filright\fontsize{60}{60}\selectfont\color{halfgray}\thechapter}
  {0pt}
  {\filright\fontsize{22}{24}\selectfont\itshape\bfseries\color{hulkblue}}
\titlespacing*{\chapter}{0pt}{-30pt}{40pt}

\titleformat{\section}
  {\normalfont\Large\bfseries}
  {\color{halfgray}\thesection\hspace{0.5em}}
  {0pt}{}

\titleformat{\subsection}
  {\normalfont\large\itshape\bfseries}
  {\thesubsection\hspace{0.5em}}
  {0pt}{}

% ── Encabezados y pie de página ─────────────────────────────────────────────
\usepackage{fancyhdr}
\pagestyle{fancy}
\fancyhf{}
\renewcommand{\headrulewidth}{0pt}
\renewcommand{\footrulewidth}{0pt}
\fancyhead[LE]{\small\color{halfgray}\thepage\quad\textit{\nouppercase{\leftmark}}}
\fancyhead[RO]{\small\color{halfgray}\textit{\nouppercase{\rightmark}}\quad\thepage}
\renewcommand{\chaptermark}[1]{\markboth{#1}{}}
\renewcommand{\sectionmark}[1]{\markright{\thesection\ #1}{}}

\fancypagestyle{plain}{%
  \fancyhf{}\renewcommand{\headrulewidth}{0pt}
  \fancyfoot[C]{\small\color{halfgray}\thepage}
}

% ── Indice: estilo clásico al modo Knuth/Bringhurst ───────────────────────
\usepackage{tocloft}
% Capítulos: versalitas, sin puntos guía, numeración alineada
\renewcommand{\cftchapfont}{\normalfont\bfseries}
\renewcommand{\cftchappagefont}{\normalfont\bfseries}
\renewcommand{\cftchapleader}{\hfill}
\renewcommand{\cftchapaftersnum}{\quad}
\renewcommand{\cftchapaftersnumb}{}
\setlength{\cftchapnumwidth}{2.2em}
\setlength{\cftbeforechapskip}{0.8em}
% Secciones: cuerpo normal, con puntos guía discretos
\renewcommand{\cftsecfont}{\normalfont}
\renewcommand{\cftsecpagefont}{\normalfont}
\renewcommand{\cftsecleader}{\cftdotfill{\cftdotsep}}
\renewcommand{\cftdotsep}{4.5}
\setlength{\cftsecindent}{2.2em}
\setlength{\cftsecnumwidth}{2.6em}
% Subsecciones
\renewcommand{\cftsubsecfont}{\normalfont\itshape}
\renewcommand{\cftsubsecpagefont}{\normalfont}
\renewcommand{\cftsubsecleader}{\cftdotfill{\cftdotsep}}
\setlength{\cftsubsecindent}{4.5em}
\setlength{\cftsubsecnumwidth}{3em}
% Título "Índice general"
\renewcommand{\contentsname}{Índice general}

% ── Hiperlinks ──────────────────────────────────────────────────────────────
\usepackage[
  pdfusetitle,
  pdfencoding=auto,
  hidelinks,
  colorlinks=false,
  linktoc=all,
  bookmarksopen=true,
  bookmarksdepth=2
]{hyperref}

% ── Cajas, listas, tablas ───────────────────────────────────────────────────
\usepackage{enumitem}
\setlist{topsep=2pt,itemsep=2pt,parsep=0pt}
\usepackage{booktabs}
\usepackage{tabularx}
\usepackage{multirow}
\usepackage{makecell}
\usepackage{array}
\usepackage{float}
\usepackage{caption}
\captionsetup{font=small,labelfont={bf,sc},labelsep=period}
\usepackage{subcaption}

% ── Cajas estilizadas ───────────────────────────────────────────────────────
\usepackage[most]{tcolorbox}
\tcbset{
  defstyle/.style={
    enhanced, breakable,
    colback=hulkblue!4, colframe=hulkblue!50!black,
    coltitle=white, fonttitle=\bfseries\scshape,
    boxrule=0.5pt, arc=2pt,
    left=10pt, right=10pt, top=6pt, bottom=6pt,
    title filled, attach boxed title to top left={xshift=10pt,yshift=-7pt},
    boxed title style={colback=hulkblue!80!black, sharp corners, boxrule=0pt}
  },
  notebox/.style={
    enhanced, breakable,
    colback=halfgray!8, colframe=halfgray!60,
    boxrule=0.5pt, arc=1pt,
    left=10pt, right=10pt, top=4pt, bottom=4pt
  }
}

\newtcolorbox{definition}[1][]{defstyle, title={Definición}, #1}
\newtcolorbox{designnote}[1][]{notebox, title={\textsc{Nota de diseño}}, #1, attach boxed title to top left={xshift=8pt,yshift=-6pt}, coltitle=halfgray!80!black, boxed title style={colback=halfgray!10, sharp corners, boxrule=0pt}}

% ── Código: estilos para HULK, Rust, LLVM IR ────────────────────────────────
\usepackage{listings}

\lstdefinestyle{base}{
  inputencoding=utf8,
  extendedchars=true,
  literate=
    {á}{{\'a}}1 {é}{{\'e}}1 {í}{{\'i}}1 {ó}{{\'o}}1 {ú}{{\'u}}1
    {Á}{{\'A}}1 {É}{{\'E}}1 {Í}{{\'I}}1 {Ó}{{\'O}}1 {Ú}{{\'U}}1
    {ñ}{{\~n}}1 {Ñ}{{\~N}}1 {ü}{{\"u}}1 {Ü}{{\"U}}1
    {¿}{{?`}}1 {¡}{{!`}}1,
  basicstyle=\ttfamily\footnotesize,
  keywordstyle=\color{codekeyword}\bfseries,
  stringstyle=\color{codestring},
  commentstyle=\color{codecomment}\itshape,
  numberstyle=\tiny\color{halfgray},
  backgroundcolor=\color{codebg},
  frame=single,
  framerule=0.4pt,
  rulecolor=\color{codeborder},
  numbers=left, numbersep=10pt,
  breaklines=true, breakatwhitespace=true,
  showstringspaces=false,
  captionpos=b, abovecaptionskip=8pt,
  belowcaptionskip=2pt,
  tabsize=2,
  xleftmargin=6pt, xrightmargin=2pt,
  framexleftmargin=4pt
}

\lstdefinelanguage{hulk}{
  morekeywords={let,in,if,elif,else,while,for,function,type,new,self,inherits,
    protocol,extends,is,as,true,false,null,base},
  morekeywords=[2]{Number,Boolean,String,Object,Null,Vector,Range,Iterable,Enumerable},
  sensitive=true,
  morecomment=[l]{//},
  morecomment=[s]{/*}{*/},
  morestring=[b]"
}
\lstdefinestyle{hulkst}{style=base,language=hulk,
  keywordstyle=[2]\color{codetype}\itshape}

\lstdefinelanguage{rustlang}{
  morekeywords={fn,let,mut,pub,struct,enum,impl,match,Ok,Err,Some,None,
    use,mod,crate,super,self,return,if,else,while,for,in,loop,trait,
    where,as,unsafe,extern,static,const,ref,move},
  morekeywords=[2]{Self,String,Vec,Box,Option,Result,bool,u32,u64,i32,
    i64,f32,f64,char,usize,isize},
  sensitive=true,
  morecomment=[l]{//},
  morecomment=[s]{/*}{*/},
  morestring=[b]"
}
\lstdefinestyle{rustst}{style=base,language=rustlang,
  keywordstyle=[2]\color{codetype}}

\lstdefinelanguage{llvm}{
  morekeywords={define,declare,alloca,store,load,call,ret,br,add,sub,mul,
    fdiv,fadd,fmul,fsub,fptoui,fptosi,uitofp,getelementptr,icmp,fcmp,phi,
    unreachable,select,and,or,xor,zext,sext,trunc,bitcast,inttoptr,ptrtoint,
    global,constant,type,private,internal,external,linkonce,common},
  morekeywords=[2]{i1,i8,i16,i32,i64,double,float,ptr,void,label},
  sensitive=true,
  morecomment=[l]{;},
  morestring=[b]"
}
\lstdefinestyle{llvmst}{style=base,language=llvm,
  keywordstyle=[2]\color{codetype}}

\lstdefinestyle{bnf}{style=base,
  keywordstyle=\color{codekeyword},
  morekeywords={},
  numbers=none}

% Etiqueta "Listado X.Y" en lugar de "Listing"
\renewcommand{\lstlistingname}{Listado}
\renewcommand{\lstlistlistingname}{Indice de listados}

% ── Diagramas ───────────────────────────────────────────────────────────────
\usepackage{tikz}
\usetikzlibrary{
  arrows.meta, automata, positioning, shapes.geometric,
  shapes.misc, calc, fit, backgrounds, decorations.pathreplacing,
  matrix, trees, chains, shapes.multipart, shadows
}

% ── Referencias cruzadas inteligentes ───────────────────────────────────────
\usepackage[spanish,capitalise,nameinlink]{cleveref}

% ── Letras capitulares y citaciones ─────────────────────────────────────────
\usepackage{lettrine}
\setcounter{secnumdepth}{3}
\setcounter{tocdepth}{2}

% ── Comandos personalizados ─────────────────────────────────────────────────
\newcommand{\hulk}{\textsc{Hulk}}
\newcommand{\code}[1]{\texttt{#1}}

% ── Mejor justificación: microtype + tolerancia ─────────────────────────────
\usepackage[activate={true,nocompatibility},final,tracking=true,kerning=true,spacing=true,factor=1100,stretch=10,shrink=10]{microtype}
\setlength{\emergencystretch}{3em}
\tolerance=1500
\hbadness=2000
\newcommand{\rust}{\textsc{Rust}}
\newcommand{\llvm}{\textsc{Llvm}}
\newcommand{\inkwell}{\textit{inkwell}}
\newcommand{\todochk}{\textsuperscript{\color{webgreen}\ding{52}}}
\usepackage{pifont}
\newcommand{\done}{{\color{webgreen}\ding{52}}}
\newcommand{\miss}{{\color{webbrown}\ding{56}}}
\newcommand{\halfdone}{{\color{orange}\ding{109}}}

% ── Sin sangria de parrafo, separación vertical en su lugar ──────────────────
\usepackage{parskip}
\setlength{\parskip}{6pt plus 2pt minus 1pt}

% ════════════════════════════════════════════════════════════════════════════
\begin{document}
% ════════════════════════════════════════════════════════════════════════════

% ────────────────────────────────────────────────────────────────────────────
%  Portada
% ────────────────────────────────────────────────────────────────────────────
\begin{titlepage}
\thispagestyle{empty}
\centering
\vspace*{1cm}

{\scshape\large Universidad de La Habana\\[0.3em]
Facultad de Matemática y Computación\par}

\vspace{0.5cm}
{\color{halfgray}\rule{0.7\textwidth}{0.4pt}}
\vspace{2.5cm}

{\fontsize{36}{40}\selectfont\itshape Un compilador para \par}
\vspace{0.6cm}
{\fontsize{72}{72}\selectfont\bfseries\color{hulkblue} HULK \par}
\vspace{1.0cm}
{\Large\itshape Diseño, implementación y extensión de un \\[0.2em]
lenguaje de programación orientado a objetos\par}

\vspace{2.5cm}

{\color{halfgray}\rule{0.4\textwidth}{0.4pt}}
\vspace{0.8cm}

{\Large
\textsc{Reporte de Proyecto Final}\\[0.5em]
\textit{Asignaturas de Compilación y Lenguajes de Programación}\par}

\vspace{1.5cm}

\begin{tabular}{rl}
\textsc{Curso}      & 2025--2026 \\[2pt]
\textsc{Carrera}    & Ciencias de la Computación \\[2pt]
\textsc{Año}        & Tercer año \\[2pt]
\textsc{Backend}    & \llvm{} 17 (código nativo x86-64) \\[2pt]
\textsc{Stack}      & \rust{} 1.77+ \\[2pt]
\textsc{Extensión}  & Vectores y expresiones generadoras \\
\end{tabular}

\vfill
{\color{halfgray}\rule{0.7\textwidth}{0.4pt}}
\vspace{0.5em}
{\large \today \par}
\end{titlepage}

% ────────────────────────────────────────────────────────────────────────────
%  Resumen
% ────────────────────────────────────────────────────────────────────────────
\thispagestyle{plain}
\chapter*{Resumen}
\addcontentsline{toc}{chapter}{Resumen}
\markboth{Resumen}{Resumen}

\lettrine[lines=2,findent=2pt,nindent=0pt]{S}{e} presenta un compilador \textit{ahead
of time} para \hulk{} (\textit{Havana University Language for Kompilers}), un
lenguaje multiparadigma estáticamente tipado, orientado a expresiones y con un
sistema de objetos basado en herencia simple y protocolos estructurales. El
compilador está escrito íntegramente en \rust{} (edición~2024) y produce
ejecutables nativos x86-64 mediante \llvm{}~17, accedido a través de la biblioteca
\inkwell{}. El \textit{frontend} (analizadores léxico, sintáctico y semántico) y
el \textit{backend} (generación de \llvm{}~IR, optimización y enlace con un
\textit{runtime} en~C) se implementan sin recurrir a generadores externos de
\textit{lexers} ni \textit{parsers}: el analizador léxico se construye sobre
autómatas finitos no deterministas obtenidos por la construcción clásica de
Thompson, y el analizador sintáctico es un motor LALR(1) cuyas tablas se generan
en tiempo de inicialización a partir de una definición declarativa de la gramática
de \hulk{}.

La implementación cubre todas las características nucleares del lenguaje
---expresiones aritméticas y de cadena, control de flujo expresivo, funciones de
primer orden, tipos con herencia simple, protocolos con conformidad estructural y
los operadores de prueba y conversión de tipos \code{is}/\code{as}---, e incorpora
como aporte propio un \textbf{sistema de vectores con expresiones generadoras}.
Esta extensión añade un constructor de tipos paramétrico $\code{Vector<T>}$ al
sistema de tipos, tres nuevas formas sintácticas (literales, generadores e
indexación) y una semántica operacional ansiosa con verificación de límites en
tiempo de ejecución. Su diseño está inscrito en la familia clásica de las
\textit{list comprehensions} de Haskell~\cite{peytonjones2003} y Python, las
\textit{for-comprehensions} de Scala~\cite{odersky2008} y los \textit{ranges and
views} de C++20, y se compara con cada una de ellas para justificar las
decisiones de diseño tomadas.

El cuerpo del documento analiza con detalle las cinco fases del compilador,
expone las decisiones técnicas relevantes junto a sus alternativas descartadas,
formaliza la sintaxis y la semántica del lenguaje implementado mediante reglas
de inferencia tipográfica y semántica operacional de gran paso, y reporta los
resultados de validación obtenidos al ejecutar la suite de pruebas sobre el
binario producido.

% ────────────────────────────────────────────────────────────────────────────
%  Indice general
% ────────────────────────────────────────────────────────────────────────────
\tableofcontents

\cleardoublepage

% ════════════════════════════════════════════════════════════════════════════
\chapter{Introducción}
% ════════════════════════════════════════════════════════════════════════════

\lettrine[lines=2,findent=2pt,nindent=0pt]{L}{a} traducción de programas escritos en
un lenguaje de alto nivel a código nativo ejecutable es uno de los procesos
fundamentales de la informática contemporánea y la actividad sobre la que se
sostienen prácticamente todas las herramientas modernas de desarrollo de
\textit{software}. Tras cinco décadas de evolución la disciplina ha consolidado un
\textit{pipeline} de fases bien diferenciadas ---análisis léxico, sintáctico y
semántico, generación de código intermedio, optimización y emisión de código
máquina--- cuyos fundamentos teóricos están extensamente documentados en la
literatura clásica~\cite{aho2006,cooper2011,appel1998}, pero cuya implementación
concreta para un lenguaje real exige resolver una gran cantidad de problemas
prácticos no triviales: representación de tipos algebraicos en memoria, despacho
de métodos en presencia de polimorfismo, propagación de tipos inferidos a través
de fases independientes, integración con \textit{backends} maduros como
\llvm{}~\cite{lattner2004}, manejo de errores con localización precisa, y un largo
etcétera.

Este reporte describe \textbf{el diseño y la construcción de un compilador
completo} para \hulk{} (\textit{Havana University Language for Kompilers}), un
lenguaje multiparadigma de propósito didáctico cuya especificación combina
elementos del paradigma funcional, imperativo y orientado a objetos. El compilador toma como entrada un fichero fuente \texttt{.hulk} y
produce un ejecutable nativo Linux x86-64 que el sistema operativo puede cargar y
ejecutar sin más dependencias externas que \code{libc} y \code{libm}. Toda la
implementación está escrita en \rust{} salvo una pequeña biblioteca de soporte en
tiempo de ejecución escrita en~C, y utiliza \llvm{}~17 como infraestructura para
la optimización y la emisión de código máquina.

\section{El lenguaje \hulk{}}

\hulk{} es un lenguaje estáticamente tipado, orientado a expresiones, con un
sistema de objetos basado en herencia simple y protocolos estructurales. Tres
rasgos lo definen como objetivo de compilación particularmente interesante.

\paragraph{Expresividad funcional.}
En \hulk{} no existe la noción de \textit{sentencia}: todas las construcciones
del lenguaje son expresiones que producen un valor. Un \code{if-elif-else} no
ejecuta ramas distintas con efecto de lado; \emph{retorna} el valor de la rama
elegida. Un bucle \code{while} retorna el valor de su última iteración. Un bloque
$\{\,e_1;\, e_2;\, \dots;\, e_n\,\}$ retorna $e_n$. Esta uniformidad obliga al
verificador de tipos a calcular el tipo de cada construcción compuesta como el
\emph{ancestro común más cercano} de los tipos de sus subexpresiones, y obliga al
generador de código a sintetizar nodos $\phi$ en cada punto de unión del grafo de
flujo.

\paragraph{Tipado nominal con protocolos estructurales.}
La jerarquía de tipos definidos por el usuario sigue el modelo nominal clásico
(dos tipos son iguales si tienen el mismo nombre), con herencia simple y
sobrescritura de métodos. Sobre esa jerarquía se superpone un segundo mecanismo
de polimorfismo basado en \textit{protocolos}: tipos abstractos que declaran
firmas de métodos sin cuerpo y que un tipo concreto \emph{satisface
implícitamente} con sólo declarar métodos compatibles ---sin necesidad de
mención explícita de la relación---. La conformidad estructural se rige por las
reglas estándar de varianza (covarianza en retornos, contravarianza en
parámetros)~\cite{pierce2002}, lo que la convierte en una construcción rica desde
el punto de vista del sistema de tipos.

\paragraph{Inferencia local de tipos.}
Las anotaciones de tipo son opcionales en parámetros, atributos e
inicializadores de \code{let}. Cuando se omiten, el verificador de tipos debe
inferir el tipo correcto mediante un análisis bidireccional: las firmas conocidas
restringen los tipos hacia abajo, los usos restringen los tipos hacia arriba.
La inferencia debe operar también en presencia de \emph{recursión mutua} entre
funciones y entre tipos, lo que obliga a separar la recolección de declaraciones
de la verificación de los cuerpos.

\section{Principios de diseño del compilador}

El compilador descrito aquí está construido en torno a cinco principios de diseño
explícitos, todos ellos consecuencia de un mismo compromiso de fondo: priorizar
la transparencia de la maquinaria interna sobre la conveniencia inmediata. Las
herramientas modernas de generación de \textit{frontends} (LALRPOP, ANTLR,
Bison) y de \textit{backends} (Cranelift, GCC plugins) permiten obtener un
compilador funcional en muy pocas líneas, pero a costa de ocultar precisamente
los algoritmos que el documento se propone analizar.

\begin{description}[style=nextline,leftmargin=2.2em]
  \item[D1. Transparencia algorítmica.]
    Cada fase implementa de forma explícita los algoritmos canónicos asociados
    a ella: construcción de Thompson en el \textit{lexer}, cálculo de FIRST y
    FOLLOW por punto fijo y construcción canónica de LR(1) con fusión LALR(1)
    en el \textit{parser}, propagación bidireccional de tipos con dos pasadas
    en el analizador semántico, \textit{lowering} estructurado a SSA en el
    generador de código. Ninguna fase delega su núcleo algorítmico en una
    biblioteca externa.
  \item[D2. Compilación a código nativo, no interpretación.]
    El compilador es un traductor estricto: produce un binario ELF~x86-64
    enlazable, no un interpretador del AST ni una máquina virtual ad-hoc. La
    intención es validar la implementación contra la métrica más exigente
    posible ---tiempos de ejecución comparables con el código~C equivalente---,
    no contra una semántica abstracta.
  \item[D3. Seguridad de memoria del compilador mismo.]
    El compilador maneja estructuras de datos complejas con múltiples
    referencias cruzadas (tablas de símbolos jerárquicas, grafos de tipos,
    tablas LALR con cientos de estados). Se utilizó \rust{} como lenguaje
    anfitrión para eliminar de raíz una clase entera de errores que
    históricamente han plagado los compiladores escritos en C++:
    \textit{dangling pointers}, \textit{double frees}, lecturas de memoria
    invalida. El \textit{wrapper} \inkwell{}~\cite{inkwell} aporta la misma
    garantía sobre el uso de la API~C de \llvm{}.
  \item[D4. Reproducibilidad y \textit{debuggability}.]
    Toda estructura de datos significativa ---la gramática, la tabla LALR, el
    autómata, la jerarquía de tipos, el \llvm{}~IR generado--- es inspeccionable
    en tiempo de ejecución mediante volcados textuales legibles. La gramática
    se describe como una estructura de datos en memoria en lugar de un
    \textit{DSL} externo, lo que permite imprimir cada producción con su
    identificador y depurar conflictos sin recompilar.
  \item[D5. Extensibilidad sintáctica y semántica.]
    El sistema de tipos, el AST y el generador de código fueron diseñados con
    anticipación a la introducción de extensiones. La extensión central
    documentada en este reporte ---vectores con expresiones generadoras---
    añade tres formas sintácticas y un nuevo constructor de tipos
    paramétrico, y demuestra empíricamente que la organización modular es
    capaz de absorber esa adición sin reescrituras significativas.
\end{description}


% ════════════════════════════════════════════════════════════════════════════
\chapter{Arquitectura del compilador}
\label{cap:arq}
% ════════════════════════════════════════════════════════════════════════════

\lettrine[lines=2,findent=2pt,nindent=0pt]{E}{l} compilador adopta la organización
en fases secuenciales que es habitual en la disciplina desde los primeros
compiladores prácticos~\cite{wirth1971,aho2006}: el programa fuente atraviesa
una sucesión de módulos, cada uno de los cuales lo transforma en una
representación interna distinta hasta llegar al binario ejecutable. Este
capítulo describe la arquitectura concreta del compilador implementado;
los capítulos sucesivos entran en el detalle de cada fase.

\section{El pipeline de compilación}

La compilación es siempre una traducción \textit{ahead of time}: el binario
\code{hulk} recibe como argumento la ruta de un fichero \texttt{.hulk}, lo
procesa por completo, escribe un fichero objeto en un directorio temporal
y lo enlaza con la biblioteca de \textit{runtime} mediante \code{gcc}, dejando
en el directorio actual un ejecutable llamado \code{output}. El proceso es
estrictamente secuencial: cada fase se ejecuta una sola vez, en orden, y
sus errores se reportan según el código de salida fijado por el contrato de
interfaz (1 para errores léxicos, 2 para sintácticos, 3 para semánticos).

La \cref{fig:pipeline} representa el pipeline tal como está implementado en
\code{src/main.rs}. Las cajas verticales en línea continua corresponden a las
fases que el compilador ejecuta en su orden de invocación; las cajas
laterales en línea punteada representan las estructuras de datos que cada
fase produce y la fase siguiente consume.

\begin{figure}[H]
\centering
\begin{tikzpicture}[
  font=\small,
  every node/.style={font=\small},
  phase/.style = {
    rectangle, draw=hulkblue!80!black, line width=0.8pt,
    top color=hulkblue!2, bottom color=hulkblue!18,
    rounded corners=3pt,
    minimum width=4.2cm, minimum height=1.0cm, align=center,
    drop shadow={shadow xshift=0.6pt, shadow yshift=-0.6pt, opacity=0.25}
  },
  back/.style = {
    rectangle, draw=hulkaccent!80!black, line width=0.8pt,
    top color=hulkaccent!2, bottom color=hulkaccent!18,
    rounded corners=3pt,
    minimum width=4.2cm, minimum height=1.0cm, align=center,
    drop shadow={shadow xshift=0.6pt, shadow yshift=-0.6pt, opacity=0.25}
  },
  io/.style = {
    rectangle, draw=halfgray, line width=0.5pt, dashed,
    fill=white,
    minimum width=4.2cm, minimum height=0.7cm, align=center,
    font=\footnotesize\itshape
  },
  ir/.style = {
    rectangle, draw=halfgray, line width=0.4pt, dashed,
    fill=halfgray!4,
    minimum width=3.4cm, minimum height=0.6cm, align=center,
    font=\footnotesize\itshape
  },
  arr/.style  = {-{Stealth[length=2.6mm,width=2.2mm]}, line width=0.55pt, color=halfgray!60!black},
  sidearr/.style = {{Stealth[length=2.6mm,width=2.2mm]}-, line width=0.4pt, dashed, color=halfgray!60!black}
]
  % main column
  \node[io]    (src)               {\texttt{programa.hulk}};
  \node[phase] (lex)  [below=8mm of src]  {Análisis léxico\\\footnotesize\code{lexer::Lexer}};
  \node[phase] (par)  [below=8mm of lex]  {Análisis sintáctico\\\footnotesize\code{parser::ParserDriver}};
  \node[phase] (sem)  [below=8mm of par]  {Análisis semántico\\\footnotesize\code{semantic::TypeChecker}};
  \node[back]  (cg)   [below=8mm of sem]  {Generación de código\\\footnotesize\code{codegen::CodegenContext}};
  \node[back]  (opt)  [below=8mm of cg]   {Optimización \llvm{}\\\footnotesize\code{mem2reg, instcombine, \dots}};
  \node[back]  (emit) [below=8mm of opt]  {Emisión de código objeto\\\footnotesize\code{TargetMachine::write\_to\_file}};
  \node[back]  (link) [below=8mm of emit] {Enlace con \textit{runtime}\\\footnotesize\code{gcc} \texttt{+} \code{hulk\_runtime.a}};
  \node[io]    (bin)  [below=8mm of link] {\texttt{./output} (ELF x86-64)};

  % main arrows
  \foreach \a/\b in {src/lex, lex/par, par/sem, sem/cg, cg/opt, opt/emit, emit/link, link/bin}
      \draw[arr] (\a) -- (\b);

  % lateral IR labels (right side)
  \node[ir] (toks)  [right=12mm of lex.east] {flujo de tokens};
  \node[ir] (ast1)  [right=12mm of par.east] {AST};
  \node[ir] (ast2)  [right=12mm of sem.east, align=center] {AST \texttt{+}\\tipos inferidos};
  \node[ir] (ir1)   [right=12mm of cg.east]  {módulo \llvm{}~IR};
  \node[ir] (ir2)   [right=12mm of opt.east] {IR optimizado};
  \node[ir] (obj)   [right=12mm of emit.east]{archivo \texttt{.o}};

  \foreach \p/\d in {lex/toks, par/ast1, sem/ast2, cg/ir1, opt/ir2, emit/obj}
      \draw[sidearr] (\d.west) -- (\p.east);

  % brace: frontend / backend zone markers (left side)
  \draw[decorate, decoration={brace, amplitude=6pt}, color=hulkblue!70!black, line width=0.6pt]
      ([xshift=-6mm]lex.north west) -- ([xshift=-6mm]sem.south west)
      node[midway, left=8pt, font=\footnotesize\scshape, color=hulkblue!70!black, rotate=90, anchor=south]
      {Frontend};
  \draw[decorate, decoration={brace, amplitude=6pt}, color=hulkaccent!70!black, line width=0.6pt]
      ([xshift=-6mm]cg.north west) -- ([xshift=-6mm]link.south west)
      node[midway, left=8pt, font=\footnotesize\scshape, color=hulkaccent!70!black, rotate=90, anchor=south]
      {Backend};
\end{tikzpicture}
\caption[Pipeline de compilación]{Pipeline de compilación implementado en
  \code{src/main.rs}. La columna central contiene las fases ejecutadas en
  orden; cada caja indica el módulo de Rust que realiza la transformación.
  Los recuadros laterales en línea punteada son las estructuras de datos que
  cada fase produce y la fase siguiente consume. Las llaves a la izquierda
  agrupan visualmente el \textit{frontend} (azul) y el \textit{backend}
  (rojo); la frontera entre ambos coincide con la producción del módulo de
  \llvm{}~IR.}
\label{fig:pipeline}
\end{figure}

\section{Las cuatro fases del \textit{frontend}}

El \textit{frontend} comprende las tres primeras fases del compilador y se
encarga de transformar el texto fuente en un AST anotado con tipos. La
implementación reparte esta responsabilidad entre tres módulos de Rust con
fronteras estrictas: \code{lexer}, \code{parser} y \code{semantic}. La
\cref{fig:pipeline} muestra cada uno alimentando al siguiente sin
retroalimentación.

\paragraph{Lexer.} El módulo \code{lexer} reconoce tokens a partir del flujo
de caracteres del fichero fuente. Su núcleo es un autómata finito no
determinista construido por el algoritmo de Thompson a partir de las
expresiones regulares declaradas para cada categoría léxica. Los autómatas
de todos los patrones se combinan en un \emph{MasterNFA}, del que el lexer
toma decisiones por la regla del máximo prefijo con desempate por prioridad.
El resultado es un vector de tokens entregado por completo al
\textit{parser}; los tokens que encubren errores léxicos se marcan como
\code{ERROR} sin interrumpir el reconocimiento, lo que permite reportar
varios errores léxicos en una misma compilación.

\paragraph{Parser.} El módulo \code{parser} consume la secuencia de tokens y
construye el árbol sintáctico abstracto. Internamente se subdivide en cuatro
submódulos: \code{grammar} define la gramática como una colección de
producciones declarativas, \code{lalr} construye las tablas LALR(1) de
\textsc{Action} y \textsc{Goto} en tiempo de inicialización a partir de esa
gramática, \code{engine} ejecuta el algoritmo \textit{shift-reduce} sobre las
tablas precalculadas, y \code{ast} define la jerarquía de nodos del árbol
sintáctico abstracto. La construcción de las tablas ---algoritmo canónico de
DeRemer~\cite{deremer1969}--- ocurre una sola vez al iniciarse el
compilador.

\paragraph{Semantic.} El módulo \code{semantic} verifica el AST y lo enriquece
con información de tipos. Se compone de un submódulo de \emph{sistema de
tipos} ---definiciones del enum \code{HulkType} y de la estructura
\code{TypeHierarchy}---, una tabla de símbolos jerárquica con
\textit{scoping} estructurado, y un \emph{type checker} que recorre el AST
en dos pasadas: la primera registra las declaraciones de tipos, protocolos
y funciones; la segunda visita los cuerpos, infiere los tipos ausentes y
escribe los resultados en un mapa indexado por el identificador único de
cada expresión. Al concluir, el análisis retorna una estructura con tres
componentes: la jerarquía de tipos resuelta, las firmas de las funciones, y
los tipos inferidos por expresión.

\section{Las fases del \textit{backend}}

El \textit{backend} ---módulo \code{codegen}--- traduce el AST tipado a un
módulo de \llvm{}~IR. La generación se orquesta en \code{codegen::jit}, que
contiene los dos puntos de entrada del módulo:
\code{execute\_program\_jit} para ejecución \textit{just-in-time} usada en
los tests, y \code{compile\_to\_binary} para la compilación AOT que produce
el binario final. Ambas comparten el mismo flujo interno.

\paragraph{Lowering.} La estructura \code{CodegenContext} encapsula el estado
de la generación: el contexto, el módulo y el constructor de \llvm{}, una
pila de tablas de símbolos con \textit{scoping} léxico, el registro de
\textit{layouts} de structs y \textit{vtables} por tipo, y los mapas de
tipos heredados del analizador semántico. La operación \code{visit\_program}
construye el módulo \llvm{} en tres pasos sucesivos: declaración de las
funciones externas del \textit{runtime}, predeclaración de todas las
funciones del programa (lo que habilita la recursión mutua sin reordenar el
AST), y construcción de los \textit{layouts} de tipos con asignación de
\textit{type tags} en orden DFS pre-orden. Sobre ese cimiento se compila el
cuerpo de cada declaración (\code{lower\_decl}) y, finalmente, la expresión
de entrada del programa (\code{lower\_expr}), envuelta en una función
\code{\_\_hulk\_entry()} accesible desde la convención de llamada de~C.

\paragraph{Optimización.} El módulo \llvm{} producido se somete a un
\textit{pipeline} fijo de transformaciones aplicadas con el \textit{new pass
manager} de \llvm{}~17, declarado en \code{codegen::jit} como la cadena
\code{"mem2reg, instcombine, reassociate, simplifycfg"}. \code{mem2reg}
es el habilitador fundamental: promueve los patrones \code{alloca/store/load}
emitidos por el \textit{lowering} a registros virtuales SSA, lo que abre la
puerta a las simplificaciones algebraicas y de control de flujo aplicadas
por los \textit{passes} subsiguientes.

\paragraph{Emisión de código objeto.} Tras la optimización, el módulo se
escribe en un fichero objeto temporal mediante
\code{TargetMachine::write\_to\_file}, configurado para la \textit{triple}
\textit{nativa} de la máquina sobre la que se ejecuta el compilador (en la
práctica, \code{x86\_64-unknown-linux-gnu}).

\paragraph{Enlace.} El compilador invoca finalmente \code{gcc} con tres
argumentos: el fichero objeto recién emitido, el archivo \code{hulk\_runtime.a}
con las primitivas de \textit{runtime}, y la bandera \code{-lm} para enlazar
\code{libm}. El resultado es un binario ELF~x86-64 ubicado en \code{./output}
que el sistema operativo carga y ejecuta sin más dependencias dinámicas que
\code{libc} y \code{libm}.

\section{Elección del lenguaje anfitrión: \rust{}}

Tres lenguajes constituyen alternativas razonables para implementar un
compilador moderno con \textit{backend} basado en \llvm{}: C, C++ y \rust{}.
La elección recayó sobre \rust{} por cuatro razones técnicas concretas.

\paragraph{Seguridad de memoria sin coste de \textit{garbage collection}.}
El compilador maneja estructuras de datos complejas con multiples referencias
cruzadas: tablas de símbolos jerárquicas, grafo de tipos con herencia, tabla de
parseo LALR con cientos de estados. \rust{} garantiza la ausencia de \textit{dangling
pointers}, \textit{double frees} y carreras de datos en tiempo de compilación, lo
que elimina una clase entera de \textit{bugs} históricamente difíciles de depurar
en compiladores escritos en C++.

\paragraph{Sistema de tipos algebraico.}
Los enumerados con datos asociados (\code{enum}) y la verificación exhaustiva de
patrones convierten la representación del AST en una tarea natural y segura. La
ausencia de un caso en un \code{match} se reporta como error en tiempo de
compilación, lo que evita bugs por casos faltantes.

\paragraph{Interfaz segura con \llvm{}.}
La biblioteca \inkwell{}~\cite{inkwell} encapsula la API~C de \llvm{} en abstracciones
de \rust{} que aprovechan el sistema de tipos para impedir operaciones invalidas
(p.\,ej. construir una instrucción en un bloque básico ya terminado). Esto
contrasta con el uso directo de \code{libLLVM} desde C++, donde tales errores
producen \textit{segfaults} difíciles de diagnosticar.

\paragraph{Ecosistema y herramientas.}
El gestor de paquetes \code{cargo}, la integración nativa con sistemas de pruebas y
la disponibilidad de bibliotecas de calidad de producción permiten enfocar el
esfuerzo en la lógica del compilador en lugar de en infraestructura.

% ════════════════════════════════════════════════════════════════════════════
\chapter{Análisis léxico}
\label{cap:lex}
% ════════════════════════════════════════════════════════════════════════════

\lettrine[lines=2,findent=2pt,nindent=0pt]{E}{l} primer paso de cualquier compilador
consiste en transformar el flujo de caracteres del fichero fuente en una secuencia
de \textit{tokens} con significado. Esta fase se denomina \textit{análisis léxico} o
\textit{tokenización}, y en este compilador se ha implementado mediante autómatas
finitos no deterministas.

\section{Arquitectura: lexer basado en NFA}

\subsection{Construcción de Thompson}

La construcción clásica de Thompson~\cite{thompson1968} permite traducir una expresión
regular cualquiera en un autómata finito no determinista con dos propiedades clave:

\begin{itemize}
  \item exactamente un estado inicial y exactamente un estado de aceptación;
  \item las transiciones epsilon ($\varepsilon$) habilitan la composición modular de
    subautómatas sin alterar el alfabeto reconocido.
\end{itemize}

La construcción procede inductivamente sobre la estructura de la expresión regular,
generando fragmentos NFA que se combinan según los operadores (concatenación,
alternación, clausura). El tamaño del NFA resultante es lineal en el tamaño de la
expresión regular original, lo que la hace adecuada para uso práctico.

\begin{figure}[H]
\centering
\begin{tikzpicture}[shorten >=1pt, node distance=1.6cm, auto,
  every state/.style={draw=hulkblue, fill=hulkblue!5, minimum size=0.8cm,
    font=\footnotesize}]
  \node[state, initial] (q0) {$q_0$};
  \node[state]          (q1) [right=of q0] {$q_1$};
  \node[state]          (q2) [right=of q1] {$q_2$};
  \node[state,accepting] (q3) [right=of q2] {$q_3$};
  \path[->] (q0) edge node {a} (q1)
            (q1) edge node {b} (q2)
            (q2) edge node {c} (q3);
\end{tikzpicture}
\caption{Fragmento NFA generado por Thompson para la expresión regular
  \code{abc}.}
\label{fig:nfa-abc}
\end{figure}

\subsection{Ejecución sin determinización}

El enfoque clásico convierte el NFA en un DFA antes de usarlo mediante la
\textit{construcción de subconjuntos}, lo que produce un autómata mas grande pero mas
rápido de simular. En este compilador se ha optado por \textbf{ejecutar el NFA
directamente}, calculando los subconjuntos de estados activos en cada paso. Esta
decisión se justifica por:

\begin{itemize}
  \item la simplicidad del código resultante;
  \item el tamaño modesto de los ficheros fuente \hulk{} (decenas a miles de líneas);
  \item la ausencia de la fase de determinización, que podria duplicar exponencialmente
    el número de estados en el peor caso.
\end{itemize}

\subsection{El MasterNFA y la regla del máximo prefijo}

Para que el \textit{lexer} reconozca todos los tipos de tokens a la vez, los NFA
individuales de cada patron se combinan en un \textbf{MasterNFA}: un autómata con
un nuevo estado inicial y transiciones epsilon hacia los estados iniciales de cada
NFA componente. La \cref{fig:master} ilustra el esquema.

\begin{figure}[H]
\centering
\begin{tikzpicture}[node distance=1.5cm, auto,
  state/.style={circle, draw=hulkblue, fill=hulkblue!5,
    minimum size=0.7cm, font=\scriptsize},
  acc/.style={state, double, fill=hulkblue!15}]
  \node[state] (s) {$s$};
  \node[state] (a1) [above right=0.5cm and 1.5cm of s] {$a_1$};
  \node[state] (a2) [right=1.5cm of s] {$b_1$};
  \node[state] (a3) [below right=0.5cm and 1.5cm of s] {$c_1$};
  \node[acc] (f1) [right=of a1] {$a_n$};
  \node[acc] (f2) [right=of a2] {$b_n$};
  \node[acc] (f3) [right=of a3] {$c_n$};
  \draw[->] (s) edge node[above,sloped] {$\varepsilon$} (a1)
            (s) edge node[above]        {$\varepsilon$} (a2)
            (s) edge node[below,sloped] {$\varepsilon$} (a3);
  \draw[->, dashed] (a1) -- (f1);
  \draw[->, dashed] (a2) -- (f2);
  \draw[->, dashed] (a3) -- (f3);
  \node[font=\scriptsize, right=0.6cm of f1] {keyword \texttt{if}};
  \node[font=\scriptsize, right=0.6cm of f2] {identificador};
  \node[font=\scriptsize, right=0.6cm of f3] {número};
\end{tikzpicture}
\caption{MasterNFA: estado inicial común con transiciones epsilon a los NFA
  individuales por patron.}
\label{fig:master}
\end{figure}

Cuando varios estados de aceptación son alcanzables simultaneamente, la regla del
\textit{máximo prefijo} (\textit{maximal munch}) selecciona el patron que reconoce la
cadena mas larga. En caso de empate (por ejemplo, \code{if} podria interpretarse como
\textit{keyword} o como identificador), la \textbf{prioridad} numérica asignada a
cada patron resuelve la ambigüedad: los \textit{keywords} tienen mayor prioridad que
los identificadores.

\section{Conjunto de tokens reconocidos}

\hulk{} define 43 tipos de token agrupados en cinco categorías. La
\cref{tab:tokens} resume la clasificación.

\begin{table}[H]
\centering
\caption{Conjunto de tokens del lenguaje \hulk{}.}
\label{tab:tokens}
\begin{tabularx}{\textwidth}{lX}
\toprule
\textsc{Categoria} & \textsc{Tokens} \\
\midrule
Palabras reservadas      &
  \code{let}, \code{in}, \code{if}, \code{elif}, \code{else},
  \code{while}, \code{for}, \code{function}, \code{type}, \code{new},
  \code{self}, \code{inherits}, \code{base}, \code{protocol},
  \code{extends}, \code{is}, \code{as}, \code{true}, \code{false}, \code{null} \\
Operadores aritméticos   &
  \code{+}, \code{-}, \code{*}, \code{/}, \code{\%}, \code{\^{}}, \code{**} \\
Operadores lógicos       &
  \code{\&}, \code{|}, \code{!} \\
Operadores de comparación &
  \code{==}, \code{!=}, \code{<}, \code{>}, \code{<=}, \code{>=} \\
Operadores especiales    &
  \code{@}, \code{@@}, \code{:=}, \code{+=}, \code{-=}, \code{*=},
  \code{/=}, \code{\%=}, \code{=>}, \code{->}, \code{++}, \code{--} \\
Delimitadores            &
  \code{(}, \code{)}, \code{\{}, \code{\}}, \code{[}, \code{]},
  \code{,}, \code{;}, \code{:}, \code{.} \\
Literales                &
  números en notación decimal y flotante, cadenas entre comillas dobles,
  caracteres entre comillas simples, identificadores alfanumericos \\
\bottomrule
\end{tabularx}
\end{table}

\section{Secuencias de escape: limitación conocida}

Los literales de cadena se reconocen como una secuencia entre comillas dobles que
admite cualquier carácter, salvo la propia comilla doble sin escapar. El
\textit{lexer} actual emite el lexema en bruto, comillas incluidas, y la fase
semántica retira las comillas envolventes para construir el nodo
\code{Literal::String}. En esta versión del compilador, las secuencias de escape
\code{\textbackslash n}, \code{\textbackslash t}, \code{\textbackslash"} y
\code{\textbackslash\textbackslash} no son descodificadas: cualquier
contrabarra que aparezca en el cuerpo del literal se trasmite tal cual al
generador de código y, por tanto, al \textit{runtime}. La consecuencia práctica es
que un \code{print("\textbackslash nhola")} imprime literalmente
\code{\textbackslash nhola} en lugar de un salto de línea seguido de
\code{hola}. Esta omisión está documentada como limitación conocida en el
\cref{cap:lim}; su implementación es directa pero queda fuera del calendario del
proyecto.

\section{Recuperación de errores}

Cuando el \textit{lexer} encuentra un carácter que no inicia ningun patron válido,
emite un token especial \code{ERROR} en lugar de detener la compilación. Esto permite
que el \textit{parser} continúe avanzando y eventualmente reporte multiples errores
en una misma invocación. El formato del mensaje sigue el contrato oficial de la
interfaz de evaluación:

\begin{lstlisting}[style=base,numbers=none,caption={Formato estandar de error léxico.}]
! (línea,columna) LEXICAL: unexpected character 'X'
\end{lstlisting}

% ════════════════════════════════════════════════════════════════════════════
\chapter{Análisis sintáctico}
\label{cap:par}
% ════════════════════════════════════════════════════════════════════════════

\lettrine[lines=2,findent=2pt,nindent=0pt]{U}{na} vez convertido el fuente en un flujo
de tokens, la fase de análisis sintáctico verifica que dicho flujo se ajuste a la
gramática del lenguaje y construye una representación estructurada del programa:
el \textit{arbol sintáctico abstracto} (AST).

\section{Parser LALR(1) implementado desde cero}

A pesar de la disponibilidad de generadores maduros como ANTLR, YACC, Bison o
LALRPOP en el ecosistema \rust{}, el parser de este compilador está implementado
manualmente sobre la base del principio de \emph{transparencia algorítmica}
(D1) enunciado en la introducción. Tres ventajas adicionales se derivan
de esa decisión:

\begin{itemize}
  \item \textbf{Inspeccionabilidad}: la gramática se representa como una
    estructura de datos \rust{} en lugar de procesarse desde un DSL textual
    externo. Esto permite imprimirla, recorrerla y analizar conflictos
    \textit{shift}/\textit{reduce} en tiempo de ejecución, sin pasos
    adicionales de compilación.
  \item \textbf{Cero dependencias externas}: el \textit{build} del compilador
    no requiere invocar ni instalar generadores como YACC o Bison. La
    cadena de \textit{build} se reduce a \code{cargo build}.
  \item \textbf{Control fino sobre la resolución de conflictos}: cuando la
    gramática introduce un conflicto deliberado ---como el del separador
    \code{|} en las expresiones generadoras descritas en el \cref{cap:ext}---,
    la resolución se programa directamente en el constructor de la tabla, sin
    necesidad de declarar precedencias en un metalenguaje externo.
\end{itemize}

\section{Pipeline del parser}

El modulo de análisis sintáctico se compone de cuatro capas claramente delimitadas
que se describen a continuación.

\subsection{Definición de la gramática}

La gramática de \hulk{} se específica en dos formas equivalentes:

\begin{itemize}
  \item un fichero \code{src/parser/grammar/hulk.grammar} en notación BNF extendida,
    legible para humanos y documentado con comentarios;
  \item una estructura \code{Vec<Production>} en \code{hulk\_grammar.rs} construida
    programáticamente, que es la versión efectivamente usada por el generador de
    tablas.
\end{itemize}

\subsection{Calculo de FIRST y FOLLOW}

Los conjuntos $\textsc{First}(\alpha)$ y $\textsc{Follow}(A)$ se calculan mediante
el algoritmo clásico de punto fijo. Recordamos las definiciones:

\begin{definition}[title={Definición: conjuntos FIRST y FOLLOW}]
Sea $G = (N, T, P, S)$ una gramática libre de contexto, $A \in N$ un no-terminal y
$\alpha \in (N \cup T)^*$ una forma sentencial.
\begin{align*}
  \textsc{First}(\alpha) &= \{ a \in T \mid \alpha \Rightarrow^* a\beta \}
    \cup \{ \varepsilon \mid \alpha \Rightarrow^* \varepsilon \} \\
  \textsc{Follow}(A) &= \{ a \in T \cup \{\$\} \mid S \Rightarrow^* \alpha A a \beta \}
\end{align*}
\end{definition}

El algoritmo de punto fijo itera sobre las producciones actualizando los conjuntos
hasta que en una iteración completa no cambia ningun conjunto.

\subsection{Construcción de la tabla LALR}

La construcción sigue el algoritmo estándar~\cite{aho2006,cooper2011} y se realiza en
tres pasos. Primero se construye la colección canónica de \textit{items} LR(1)
mediante un BFS desde el \textit{item} inicial
$[\,\textsf{Start} \rightarrow \bullet\,\textsf{Program}\,\$\,,\,\$\,]$ tomando
clausuras y aplicando $\textsc{Goto}(I, X)$ para cada par estado-símbolo. Para la
gramática completa de \hulk{} este paso produce del orden de 500--600 estados LR(1).
Segundo, se aplica el paso característico del algoritmo LALR(1)~\cite{deremer1969}:
dos estados con idéntico \textit{núcleo} (igual conjunto de pares producción-punto,
ignorando \textit{lookaheads}) se funden en uno solo, uniendo sus conjuntos de
\textit{lookaheads}. La fusión reduce el espacio de estados a aproximadamente
$\mathbf{263}$ estados, una reducción de cerca del $55\%$. Tercero, se rellenan las
tablas $\textsc{Action}[s,t]$ y $\textsc{Goto}[s,N]$ recorriendo los \textit{items}
de cada estado fusionado y se registran los conflictos detectados para su
inspección posterior.

\begin{table}[H]
\centering
\caption{Estadísticas del autómata de parseo construido para \hulk{}.}
\label{tab:lalr-stats}
\begin{tabularx}{0.75\textwidth}{Xr}
\toprule
\textsc{Métrica} & \textsc{Valor} \\
\midrule
Producciones declaradas                 &  $\sim 90$   \\
Estados LR(1) antes de fusión           &  $500$--$600$ \\
Estados LALR(1) tras fusión por núcleo  &  $\mathbf{263}$ \\
Entradas no triviales en \textsc{Action} &  ${\sim}1\,500$ \\
Entradas no triviales en \textsc{Goto}   &  ${\sim}800$  \\
Conflictos \textit{shift}/\textit{reduce} sin resolver &  $0$ \\
Conflictos \textit{shift}/\textit{reduce} resueltos    &  $1$ \\
\bottomrule
\end{tabularx}
\end{table}

El único conflicto de la gramática se sitúa sobre el terminal \code{|} y nace de la
interacción entre el operador OR lógico y el separador de las expresiones
generadoras introducidas por la extensión (\cref{cap:ext}). Cuando el motor de
parseo se encuentra en un estado con la pila de la forma
$[\,\textsf{[}\,\,\textsf{Expr}\,]$ y el siguiente terminal es \code{|}, el
\textit{lookahead} estándar de un símbolo no basta: \code{[Expr | \dots ]} puede
continuar tanto en $\textsf{[}\,\textsf{Expr}\,\textsf{|}\,\textsf{Expr}\,\textsf{]}$
---un vector explícito cuyo primer elemento es un OR lógico--- como en
$\textsf{[}\,\textsf{Expr}\,\textsf{|}\,\textsf{Id}\,\textsf{in}\,\textsf{Expr}\,\textsf{]}$
---un generador---. El constructor de la tabla resuelve este caso particular
introduciendo un \textit{lookahead} adicional de dos tokens: si tras el \code{|}
viene un \code{IDENTIFIER} seguido de \code{in}, se opta por la regla del generador;
en cualquier otro caso, por la regla de OR. La resolución se realiza en el momento
de la construcción de la tabla, no en tiempo de parseo, por lo que el motor sigue
siendo estrictamente LALR(1) en ejecución.

\subsection{Motor de parseo}

El motor implementa el algoritmo shift-reduce sobre las tablas precomputadas.
Mantiene una pila de estados y una pila de valores semánticos en paralelo. En cada
reducción se invoca una acción semántica que combina los valores semánticos del
lado derecho de la producción en un nodo del AST.

\section{Jerarquía de precedencias}

La especificación oficial de \hulk{} no establece la precedencia de operadores de
forma explicita. La gramática implementada codifica la precedencia estandar
matemática directamente en la jerarquía de no-terminales, evitando la necesidad de
declaraciones de precedencia separadas. El \cref{lst:prec} muestra el esquema.

\begin{lstlisting}[style=bnf, caption={Codificación de precedencias en la gramática.},
  label={lst:prec}]
Expr        -> Assignment
Assignment  -> OrExpr (':=' Assignment)?
OrExpr      -> AndExpr ('|' AndExpr)*
AndExpr     -> NotExpr ('&' NotExpr)*
NotExpr     -> '!' NotExpr | Equality
Equality    -> Comparison (('==' | '!=') Comparison)*
Comparison  -> Concat (('<'|'>'|'<='|'>=') Concat)*
Concat      -> Additive (('@' | '@@') Additive)*
Additive    -> Term (('+' | '-') Term)*
Term        -> Power (('*' | '/' | '%') Power)*
Power       -> Unary ('^' Power)?     (* asociatividad derecha *)
Unary       -> '-' Unary | Postfix
Postfix     -> Atom ('.' id '(' args ')' | '.' id | '[' Expr ']')*
Atom        -> literal | id | '(' Expr ')' | New | If | Let | ...
\end{lstlisting}

La asociatividad derecha de la potencia se obtiene mediante la recursión a la
derecha en la regla \code{Power}: la expresión $2 \,\widehat{}\, 3 \,\widehat{}\, 2$
se parsea como $2 \,\widehat{}\, (3 \,\widehat{}\, 2) = 2^9 = 512$, no como
$(2 \,\widehat{}\, 3) \,\widehat{}\, 2 = 64$.

\section{El arbol sintáctico abstracto}

El AST esta definido en \code{src/parser/ast/} mediante enumerados \rust{} con datos
asociados. El nodo raíz es \code{Program}, que contiene una lista de declaraciones
(\code{Decl}) y una expresión de entrada (\code{Expr}).

\begin{lstlisting}[style=rustst, caption={Estructura central del AST.},
  label={lst:ast}]
pub struct Program {
    pub declarations: Vec<Decl>,
    pub entry:        Expr,
}

pub struct Expr {
    pub kind: ExprKind,
    pub id:   u32,    // identificador único para la tabla de tipos
    pub span: Span,   // línea y columna para mensajes de error
}

pub enum ExprKind {
    Literal(Literal),
    Identifier { name: String },
    Binary(BinaryExpr),     Unary(UnaryExpr),
    Postfix(PostfixExpr),   Assign(AssignExpr),
    Let(LetExpr),           Block(BlockExpr),
    If(IfExpr),             While(WhileExpr),
    For(ForExpr),           Call(CallExpr),
    MethodCall(MethodCallExpr),
    Access(AccessExpr),     Index(IndexExpr),
    New(NewExpr),           Base,
    Is { expr: Box<Expr>, type_name: TypeName },
    As { expr: Box<Expr>, type_name: TypeName },
    Vector(VectorExpr),     // <-- extensión propia (cap. 8)
}
\end{lstlisting}

Cada \code{Expr} lleva un identificador único \code{id} que se utiliza como clave
en la tabla \code{expr\_types: HashMap<u32, HulkType>} que el analizador semántico
construye. El generador de código consulta esta tabla para conocer el tipo
inferido de cada subexpresión sin tener que reinferir.

% ════════════════════════════════════════════════════════════════════════════
\chapter{Análisis semántico}
\label{cap:sem}
% ════════════════════════════════════════════════════════════════════════════

\lettrine[lines=2,findent=2pt,nindent=0pt]{C}{onstruido} el AST, la siguiente fase
verifica que el programa tenga sentido a nivel de tipos, alcances y referencias.
Esta fase es responsable de detectar errores como uso de variables no declaradas,
asignaciones entre tipos incompatibles, llamadas a funciones con número o tipo
incorrecto de argumentos, y violaciones de la conformidad estructural a protocolos.

\section{Estrategia en dos pasadas}

El analizador semántico opera en dos pasadas claramente diferenciadas sobre la
lista de declaraciones del programa.

\paragraph{Pasada 1 --- Recolección de declaraciones.}
Todas las firmas de funciones, los miembros de tipos (atributos y métodos) y los
cuerpos de protocolos se registran en la estructura \code{TypeHierarchy}. Los
miembros sin anotación de tipo reciben un marcador \code{HulkType::Unknown} que
sera resuelto en la segunda pasada. Esta pasada es indispensable para soportar
\textbf{recursión mutua} entre funciones y tipos: en el momento de verificar el
cuerpo de \code{foo}, la firma de \code{bar} ya esta disponible aunque \code{bar}
este declarada después.

\paragraph{Pasada 2 --- Verificación de cuerpos y propagación de tipos.}
Cada cuerpo de función, cada inicializador de atributo y cada cuerpo de método se
verifica. El tipo inferido del cuerpo se compara con la anotación declarada si
existe; si no existe, el tipo inferido se \textbf{escribe de vuelta} en la
\code{TypeHierarchy} reemplazando el marcador \code{Unknown}.

Esta escritura es \emph{crítica} para la fase de generación de código. Si un atributo
\code{val = start} no tiene anotación de tipo, la pasada~1 almacena \code{Unknown}; la
pasada~2 infiere \code{Number} a partir del inicializador; la escritura de vuelta
hace que \code{Number} este disponible para el generador de código. Sin esta
escritura, el codegen no podria determinar el tipo \llvm{} correcto del campo.

\begin{lstlisting}[style=rustst, caption={Propagación del tipo inferido en la pasada
  semántica.}, label={lst:writeback}]
// Después de inferir el tipo del campo en la pasada 2:
if let Some(type_info) = self.types.types.get_mut(type_name) {
    if let Some(attr) = type_info.attributes.get_mut(field_name) {
        if *attr == HulkType::Unknown {
            *attr = inferred_type.clone();
        }
    }
}
\end{lstlisting}

\section{Sistema de tipos}

El sistema de tipos es \textbf{nominal} (la identidad de los tipos depende de su
nombre y no de su estructura) con \textbf{subtype polymorphism}~\cite{pierce2002}.

\subsection{Categorias de tipo}

\begin{itemize}
  \item \textbf{Primitivos}: \code{Number} (IEEE-754 \textit{double} de 64 bits),
    \code{Boolean} (un bit), \code{String} (puntero a cadena C
    \textit{null}-terminada), \code{Null} (tipo singleton).
  \item \textbf{Definidos por el usuario}: introducidos mediante \code{type} con
    herencia simple opcional. Todos heredan transitivamente de \code{Object}.
  \item \textbf{Protocolos}: tipos de interfaz que listan firmas de métodos. La
    conformidad es \textbf{estructural}: un tipo conforma a un protocolo si declara
    todos sus métodos con firmas compatibles, sin necesidad de declarar la
    relación explicitamente.
  \item \textbf{Vectores}: el tipo \code{Vector<T>} para colecciones homogeneas de
    elementos de tipo \code{T} (vease \cref{cap:ext}).
  \item \textbf{Tipos auxiliares}: \code{Unknown} (marcador en la pasada~1),
    \code{Never} (tipo de error para evitar cascadas de errores falsos).
\end{itemize}

\subsection{Relación de conformidad}

La relación de \textit{conformidad} $\preceq$ entre tipos formaliza el subtipado.
Se define inductivamente y opera tanto sobre la jerarquía nominal (tipos definidos
por el usuario) como sobre la jerarquía estructural (protocolos).

\begin{definition}[title={Definición: conformidad de tipos}]
Sean $A, B$ tipos del lenguaje \hulk{}. Se dice que $A$ \emph{conforma a} $B$,
notación $A \preceq B$, si y sólo si se cumple alguna de las siguientes condiciones:
\begin{align*}
  & A = B
  && \text{(reflexividad)} \\
  & A \text{ hereda transitivamente de } B
  && \text{(subtipado nominal)} \\
  & B \text{ es protocolo y } A \text{ lo implementa}
  && \text{(conformidad estructural)} \\
  & A = \mathtt{Null} \text{ y } B \text{ es nominal o } \mathtt{Object}
  && \text{(nulabilidad)} \\
  & A = \mathtt{Vec\langle S\rangle},\ B = \mathtt{Vec\langle T\rangle},\ S \preceq T
  && \text{(covarianza de vectores)} \\
  & B = \mathtt{Object}
  && \text{(raíz universal)}
\end{align*}
\end{definition}

La conformidad estructural a un protocolo $P$ no se reduce a una verificación
puramente sintáctica de presencia de métodos: cada firma debe ser \emph{compatible}
con la del protocolo, donde compatibilidad significa \textbf{covarianza en el tipo
de retorno y contravarianza en los parámetros}. Estas reglas son las clásicas para
preservar la seguridad de la sustitución en presencia de subtipado, debidas
originalmente a Liskov y Wing~\cite{pierce2002}.

\begin{definition}[title={Definición: compatibilidad de firmas con varianza}]
Sean $\sigma_A = (\bar p_A) \to r_A$ la firma del método declarado en el tipo
candidato y $\sigma_P = (\bar p_P) \to r_P$ la firma exigida por el protocolo.
$\sigma_A$ es compatible con $\sigma_P$ si y sólo si:
\begin{align*}
  & |\bar p_A| = |\bar p_P|
  & \text{(aridad)} \\
  & r_A \preceq r_P
  & \text{(retorno covariante)} \\
  & \forall i.\ p_P^{(i)} \preceq p_A^{(i)}
  & \text{(parámetros contravariantes)}
\end{align*}
\end{definition}

Estas reglas se implementan en el método
\code{signatures\_protocol\_compatible} de la jerarquía de tipos: la conformidad
se decide recorriendo los métodos del protocolo (y los del protocolo padre, si
lo tiene) y delegando en la relación \code{conforms} para cada par de tipos. La
búsqueda del método sube por la cadena de herencia del candidato, de modo que un
método heredado satisface igual de bien la firma exigida que uno propio. El
verificador atajará la búsqueda en cuanto detecte que el padre del candidato ya
conforma al protocolo, evitando reexaminar firmas heredadas.

\subsection{Unificación de tipos en condicionales: el ancestro común más cercano}

Una situación recurrente en \hulk{} es la determinación del tipo de una expresión
\code{if-elif-else} cuyas ramas devuelven valores de tipos relacionados pero no
idénticos. La regla aplicada es la habitual en lenguajes con subtipado: el tipo de
la expresión es el \emph{ancestro común más cercano} (\textit{lowest common
ancestor}, $\mathrm{lca}$) de los tipos de todas las ramas.

\begin{definition}[title={Definición: ancestro común más cercano}]
$\mathrm{lca}(A, B)$ se define como el tipo de menor profundidad en la jerarquía
de herencia que satisface $A \preceq T \land B \preceq T$. Para $n$ ramas se
define recursivamente:
$\mathrm{lca}(T_1, T_2, \ldots, T_n) = \mathrm{lca}(\mathrm{lca}(T_1, T_2), T_3,
\ldots, T_n)$.
\end{definition}

El algoritmo implementado es directo: se recolecta el conjunto $\mathcal{A}(A)$ de
todos los ancestros de $A$ (incluido $A$) subiendo por la cadena de
\code{TypeInfo.parent}, y se asciende desde $B$ hasta encontrar el primer tipo que
pertenece a $\mathcal{A}(A)$. La existencia de tal tipo está garantizada porque
todos los tipos heredan transitivamente de \code{Object}, que actúa como elemento
maximal del orden.

\section{Funciones y constantes predefinidas}

El compilador pre-registra el siguiente conjunto de elementos accesibles desde
cualquier alcance, sin necesidad de importación:

\begin{table}[H]
\centering
\caption{Funciones y constantes predefinidas del runtime \hulk{}.}
\label{tab:builtins}
\begin{tabularx}{\textwidth}{lXl}
\toprule
\textsc{Nombre} & \textsc{Descripción} & \textsc{Tipo} \\
\midrule
\code{print(x)}            & Imprime el valor en stdout                & \code{Object → Null} \\
\code{sqrt(x)}             & Raíz cuadrada                              & \code{Number → Number} \\
\code{sin(x)}              & Seno (radianes)                            & \code{Number → Number} \\
\code{cos(x)}              & Coseno (radianes)                          & \code{Number → Number} \\
\code{exp(x)}              & Exponencial $e^x$                          & \code{Number → Number} \\
\code{log(b, x)}           & Logaritmo en base $b$                      & \code{(Number, Number) → Number} \\
\code{rand()}              & Aleatorio uniforme en $[0,1)$              & \code{Number} \\
\code{range(s, e)}         & Rango iterable $[s, e)$                    & \code{(Number, Number) → Range} \\
\code{PI}, \code{E}        & Constantes matemáticas                     & \code{Number} \\
\code{true}, \code{false}  & Literales booleanos                        & \code{Boolean} \\
\bottomrule
\end{tabularx}
\end{table}

\section{Reporte de errores semánticos}

Los errores semánticos se acumulan en un vector \code{Vec<SemanticError>} y se
reportan todos juntos al final del análisis. Esta estrategia evita la frustración
del programador que debe corregir un error a la vez. Cada \code{SemanticError}
incluye el \code{Span} de la expresión o declaración causante, lo que produce
mensajes con posición precisa:

\begin{lstlisting}[style=base,numbers=none]
! (12,5) SEMANTIC: type 'Cat' does not implement method 'speak' of protocol 'Animal'
! (18,3) SEMANTIC: cannot assign value of type 'String' to variable of type 'Number'
! (24,9) SEMANTIC: function 'foo' expects 2 arguments, 3 given
\end{lstlisting}

% ════════════════════════════════════════════════════════════════════════════
\chapter{Generación de código}
\label{cap:cg}
% ════════════════════════════════════════════════════════════════════════════

\lettrine[lines=2,findent=2pt,nindent=0pt]{C}{ompletado} el análisis semántico, el
compilador dispone de un AST anotado con tipos y de una jerarquía de tipos
verificada. La fase de generación de código traduce este AST en un modulo de
\llvm{}~Intermediate~Representation, listo para ser optimizado y compilado a código
máquina nativo.

\section{\llvm{}~IR mediante \inkwell{}}

\llvm{}~\cite{lattner2004} es la infraestructura de compilación mas utilizada del
mundo. Provee una representación intermedia tipada y de bajo nivel, un conjunto
extensible de \textit{passes} de optimización, y \textit{backends} para una variedad
de arquitecturas (x86-64, ARM, RISC-V, WebAssembly, etc.).

La biblioteca \inkwell{} es un \textit{wrapper} \rust{} de la API~C de \llvm{} que
aprovecha el sistema de tipos del lenguaje para impedir construcciones invalidas en
tiempo de compilación. Por ejemplo, los \textit{builder} de instrucciones rechazan
operaciones sobre bloques básicos ya terminados, y los valores estan parametrizados
por el \code{Context} de \llvm{} para impedir mezclas entre contextos distintos.

El estado de la generación de código se centraliza en \code{CodegenContext}, una
estructura que agrupa:

\begin{itemize}
  \item \code{context}, \code{module}, \code{builder}: los tres objetos centrales
    de la API \llvm{}.
  \item \code{symbols}: pila de alcances para variables locales.
  \item \code{type\_registry}: layouts de structs \llvm{} y vtables por tipo \hulk{}.
  \item \code{type\_hierarchy}: jerarquía de tipos del analizador semántico,
    accesible para resolución de métodos en herencia.
  \item \code{expr\_types}: tipos anotados por el verificador, consultados durante
    la generación.
  \item \code{current\_function}, \code{self\_ptr}, \code{current\_type\_name}:
    contexto de la función o método en compilación.
\end{itemize}

\section{Representación de valores: el enum \code{CgValue}}

Dado que \hulk{} tiene tipos heterogeneos (\code{Number} es escalar, \code{String}
y los objetos son punteros), no existe un único tipo \llvm{} universal. La
representación uniforme se obtiene mediante un enum \rust{} que transporta el
valor \llvm{} junto con su categoría de tipo \hulk{}:

\begin{lstlisting}[style=rustst, caption={Representación de valores en codegen.},
  label={lst:cgvalue}]
pub enum CgValue<'ctx> {
    Number(FloatValue<'ctx>),    // LLVM: double
    Bool(IntValue<'ctx>),        // LLVM: i1
    Str(PointerValue<'ctx>),     // LLVM: ptr a cadena C
    Object(PointerValue<'ctx>),  // LLVM: ptr al struct del tipo
    Vector(PointerValue<'ctx>),  // LLVM: ptr al buffer en heap
    Null,                        // puntero nulo
    Void,                        // sin valor (efectos de lado)
}
\end{lstlisting}

Las coerciones entre variantes se realizan en \code{coerce\_arg}, que inserta las
instrucciones \llvm{} necesarias: \code{Bool} a \code{Number} usa \code{uitofp},
cualquier tipo primitivo a \code{ptr} para paso a \code{print(x: Object)} se
empaqueta mediante un \code{alloca + store}.

\section{Representación de objetos en memoria}

Todo tipo definido por el usuario se representa como un \code{struct} \llvm{} con
una distribución de campos uniforme. Los dos primeros campos son metadatos comunes
a todos los objetos; los siguientes son los campos declarados por el usuario en
orden \emph{padre-primero}.

\begin{figure}[H]
\centering
\begin{tikzpicture}[font=\small]
  \draw[fill=hulkblue!10, draw=hulkblue] (0,0) rectangle (2,0.8);
  \draw[fill=hulkblue!10, draw=hulkblue] (2,0) rectangle (4,0.8);
  \draw[fill=halfgray!10, draw=halfgray] (4,0) rectangle (6,0.8);
  \draw[fill=halfgray!10, draw=halfgray] (6,0) rectangle (8,0.8);
  \draw[fill=halfgray!10, draw=halfgray] (8,0) rectangle (10,0.8);
  \node at (1,0.4) {\code{i32 type\_tag}};
  \node at (3,0.4) {\code{ptr vtable}};
  \node at (5,0.4) {\code{padre.f1}};
  \node at (7,0.4) {\code{padre.f2}};
  \node at (9,0.4) {\code{propio.f1}};
  \node[font=\scriptsize] at (1,-0.3) {0};
  \node[font=\scriptsize] at (3,-0.3) {1};
  \node[font=\scriptsize] at (5,-0.3) {2};
  \node[font=\scriptsize] at (7,-0.3) {3};
  \node[font=\scriptsize] at (9,-0.3) {4};
  \draw[decorate, decoration={brace, mirror, amplitude=4pt}]
    (0,-0.5) -- (4,-0.5) node[midway,below,yshift=-4pt] {metadatos};
  \draw[decorate, decoration={brace, mirror, amplitude=4pt}]
    (4,-0.5) -- (10,-0.5) node[midway,below,yshift=-4pt] {campos en orden padre-primero};
\end{tikzpicture}
\caption{Layout en memoria de un objeto \hulk{} que hereda de un padre con dos
  campos.}
\label{fig:layout}
\end{figure}

El orden \textit{padre primero} es deliberado: garantiza que un puntero a un
objeto del tipo \code{Dog} (que hereda de \code{Animal}) sea \textbf{bit-compatible}
con un puntero a \code{Animal} respecto a los campos heredados. Esto evita el costo
de reinterpretación de bits o de tablas de desplazamiento que serian necesarias
con otros layouts.

\section{Tablas de métodos virtuales (vtables)}

El polimorfismo dinámico de \hulk{} se implementa mediante \textit{vtables}, la
técnica estandar utilizada por C++, Java (compilado JIT) y otros lenguajes con
herencia y polimorfismo.

Para cada tipo se crea una \textbf{vtable global} (un \code{global} \llvm{}) que
contiene punteros a los métodos del tipo, en un orden fijo: primero los métodos
heredados del padre, luego los nuevos. Cuando un tipo hace \textit{override} de un
método del padre, la posición en la vtable se conserva pero el puntero se
reemplaza por el del nuevo método.

\begin{lstlisting}[style=llvmst, caption={Vtable global para un tipo \code{Dog} que
  redefine \code{speak} y hereda \code{move}.}, label={lst:vtable}]
%VTable_Dog = type { ptr, ptr }

@vtable_Dog = global %VTable_Dog {
    ptr @__hulk_method_Dog_speak,      ; override
    ptr @__hulk_method_Animal_move     ; heredado del padre
}
\end{lstlisting}

El despacho de un método \code{x.speak()} sobre un valor \code{x} de tipo estático
\code{Animal} se compila como un acceso indirecto en tres pasos: cargar el puntero
a vtable desde el campo~1 del objeto, indexar al slot precalculado para el método
\code{speak}, e invocar el puntero a función con la convención de llamada habitual
(\code{self} en la primera posición, argumentos a continuación).

\subsection{Despacho de métodos a través de un protocolo}

El esquema de vtables es suficiente cuando el tipo estático del receptor es un tipo
nominal: en ese caso, el orden de los \textit{slots} es el mismo en todas las
clases de la jerarquía y la posición del método se calcula en tiempo de
compilación. Sin embargo, cuando el tipo estático es un protocolo \code{P}, los
\textit{slots} pueden diferir entre tipos conformantes: si \code{A} declara
\{\code{m},\code{n}\} y \code{B} declara \{\code{n},\code{m}\}, la posición de
\code{m} en cada vtable es distinta, y un acceso por índice fijo daría el método
equivocado.

La estrategia elegida en este compilador resuelve esa indeterminación con un
\textit{switch} sobre el \code{type\_tag} leído en tiempo de ejecución, generando
una rama por cada tipo conforme conocido y, en cada rama, una llamada
\emph{directa} al método correspondiente. La construcción es local a la expresión
de despacho ---no requiere tablas globales adicionales--- y se beneficia de la
optimización clásica que \llvm{} aplica a los \code{switch}: tras \code{simplifycfg}
se transforma en una \textit{jump table}, lo que asintóticamente tiene el mismo
coste que un acceso a vtable. El bloque \code{default} del switch se marca
\code{unreachable} porque el verificador de tipos garantiza que el receptor
realmente conforma al protocolo:

\begin{lstlisting}[style=llvmst,caption={Despacho a través de un protocolo
  \code{Measurable} con dos tipos conformes.}]
%tag = load i32, ptr %b

switch i32 %tag, label %proto_unreachable [
  i32 3, label %proto_case_MBox
  i32 5, label %proto_case_Cylinder
]

proto_case_MBox:
  %r1 = call f64 @__hulk_method_MBox_measure(ptr %b)
  br label %proto_merge

proto_case_Cylinder:
  %r2 = call f64 @__hulk_method_Cylinder_measure(ptr %b)
  br label %proto_merge

proto_unreachable:
  unreachable                  ; semánticamente inalcanzable

proto_merge:
  %result = phi f64 [ %r1, %proto_case_MBox ], [ %r2, %proto_case_Cylinder ]
\end{lstlisting}

La decisión entre despacho por vtable (para tipos nominales) y despacho por
\textit{switch} (para protocolos) se toma en \code{method\_dispatch} del módulo
\code{codegen::context} consultando la categoría sintáctica del tipo estático del
receptor.

\section{Etiquetas de tipo en orden DFS}

El operador \code{is} (\code{x is Dog}) es una operación frecuente en código OOP con
herencia, y se utiliza también internamente en la implementación del operador
\code{as} (\textit{downcast}). Implementarlo eficientemente requiere un invariante
no trivial sobre la asignación de etiquetas de tipo.

\paragraph{Asignación DFS pre-orden.}
Durante la compilación, los tipos reciben \textbf{etiquetas enteras en orden de
recorrido en profundidad de la jerarquía de herencia}, en pre-orden. El tipo padre
recibe su etiqueta antes que sus hijos. Como consecuencia, \textbf{todos los
subtipos directos e indirectos de un tipo $T$ forman un intervalo contiguo}
$[\textsf{min\_tag}_T, \textsf{max\_tag}_T]$ en el espacio de etiquetas.

\begin{figure}[H]
\centering
\begin{tikzpicture}[
  level distance=1.5cm, sibling distance=2.5cm,
  every node/.style={rectangle, draw=hulkblue, fill=hulkblue!5, rounded corners=2pt,
    minimum width=2cm, minimum height=0.6cm, font=\small},
  edge from parent/.style={draw=halfgray, -Stealth}]
  \node {Animal [1,5]}
    child { node {Dog [2,4]}
      child { node {Poodle [3,3]} }
      child { node {Beagle [4,4]} }
    }
    child { node {Cat [5,5]} };
  \node[draw=none, fill=none, right=2cm, align=left, font=\footnotesize] at (5,0)
    {\code{poodle is Dog}:    $2 \leq 3 \leq 4$ \done\\[2pt]
     \code{poodle is Animal}: $1 \leq 3 \leq 5$ \done\\[2pt]
     \code{cat is Dog}:       $2 \leq 5 \leq 4$ \miss};
\end{tikzpicture}
\caption[Asignación DFS de etiquetas]{Asignación DFS pre-orden de etiquetas de
  tipo. El test \code{is} se reduce a una verificación de pertenencia a un
  intervalo, computable en tiempo constante.}
\label{fig:dfs-tags}
\end{figure}

\paragraph{Implementación del operador \code{is}.}
Gracias a esta asignación, el código \llvm{} generado para \code{x is Dog} es
simplemente:

\begin{lstlisting}[style=llvmst, numbers=none]
  %tag = load i32, ptr %x
  %ge  = icmp uge i32 %tag, 2       ; min_tag(Dog)
  %le  = icmp ule i32 %tag, 4       ; max_tag(Dog)
  %res = and i1 %ge, %le
\end{lstlisting}

Esto es $O(1)$ en tiempo de ejecución, sin consultar ninguna tabla ni recorrer la
jerarquía. La técnica esta documentada en implementaciones eficientes de RTTI en
C++~\cite{bjarne1994}.

\section{El operador \code{as}: downcast con verificación}

El operador \code{x as T} realiza un \textit{downcast}: convierte un valor de tipo
base (por ejemplo \code{Animal}) al tipo derivado \code{T} (por ejemplo \code{Dog}).
En \hulk{}, esto debe ser \textit{seguro}: si el objeto realmente no es de tipo
\code{T}, la ejecución debe abortar con un mensaje de error en lugar de continuar
con un puntero malformado.

El código generado primero verifica el rango de etiquetas (igual que en \code{is}).
Si la verificación falla, llama a la función de \textit{runtime}
\code{hulk\_type\_error}, que imprime el mensaje de error y llama a \code{abort(3)}.
Si la verificación tiene exito, el puntero original se utiliza directamente: no se
requiere reinterpretación de bits porque el layout es compatible por la
organización padre-primero.

\section{La palabra clave \code{base()}}

Dentro de un método, \code{base(args)} invoca la implementación del padre. A
diferencia del despacho virtual --- que cargaria la vtable y podria producir una
llamada recursiva si el padre llama de vuelta al método del hijo ---, \code{base()}
es una \textbf{llamada directa estática}: el nombre de la función del padre se
conoce en tiempo de compilación.

El código generador busca recursivamente en la jerarquía hasta encontrar el tipo
que realmente implementa el método (no necesariamente el padre inmediato, si
varios niveles intermedios solo lo heredan), y emite una llamada directa por
nombre.

\section{Pipeline de optimización}

El IR generado por el codegen es deliberadamente sencillo y abundante en
\textit{allocas}: cada variable local se modela como una secuencia
\code{alloca + store + load}. Esta forma facilita la generación correcta pero es
ineficiente sin optimización.

Antes de la emisión de código objeto se aplica el \textit{pipeline} del nuevo
\textit{pass manager} de \llvm{}~17 con la siguiente secuencia:

\begin{table}[H]
\centering
\caption{Passes de optimización aplicados al IR.}
\label{tab:passes}
\begin{tabularx}{\textwidth}{lX}
\toprule
\textsc{Pass} & \textsc{Efecto} \\
\midrule
\code{mem2reg} & Promueve patrones \code{alloca}/\code{store}/\code{load} a
  registros SSA virtuales. Es el habilitador fundamental de todas las
  optimizaciones subsiguientes. \\
\code{instcombine} & Plegado de constantes, simplificación algebraica, eliminación
  de operaciones redundantes (por ejemplo, $x \oplus 0 = x$, $x \cdot 1 = x$). \\
\code{reassociate} & Reordena expresiones asociativas para maximizar la oportunidad
  de plegado por \code{instcombine}. \\
\code{simplifycfg} & Elimina bloques básicos muertos, fusiona bloques redundantes,
  simplifica el grafo de flujo de control. \\
\bottomrule
\end{tabularx}
\end{table}

Esta combinación produce un IR significativamente mas compacto y próximo al código
generado por compiladores maduros como \code{gcc -O2}.

% ════════════════════════════════════════════════════════════════════════════
\chapter{Runtime y compilación final}
\label{cap:rt}
% ════════════════════════════════════════════════════════════════════════════

\lettrine[lines=2,findent=2pt,nindent=0pt]{E}{l} código \llvm{} producido por el
\textit{frontend} no es autosuficiente: depende de un conjunto de funciones de
soporte para operaciones como impresión, formato de números, allocación de cadenas
y manejo de colecciones. Estas funciones forman la \textit{biblioteca de runtime},
escrita en~C y enlazada con cada ejecutable producido.

\section{La biblioteca de runtime en C}

El \textit{runtime} se compila a un archivo estático \code{hulk\_runtime.a} desde
\code{runtime/hulk\_runtime.c}. Provee las siguientes funciones, todas con enlace~C
y nombres no decorados para que sean accesibles desde el código generado.

\begin{table}[H]
\centering
\caption{Funciones de la biblioteca de runtime.}
\label{tab:runtime}
\begin{tabularx}{\textwidth}{lX}
\toprule
\textsc{Función} & \textsc{Descripción} \\
\midrule
\code{hulk\_print(ptr)}            & Imprime cadena $+$ \textit{newline} a stdout. \\
\code{hulk\_str\_from\_number(d)}  & Convierte \code{double} a C string;
  formatea enteros sin parte decimal. \\
\code{hulk\_str\_concat(a,b)}      & Concatenación para el operador \code{@}. \\
\code{hulk\_str\_concat\_space(a,b)} & Concatenación con espacio para
  \code{@@}. \\
\code{hulk\_str\_eq(a,b)}          & Igualdad de cadenas por valor. \\
\code{hulk\_str\_size(s)}          & Longitud en bytes de la cadena. \\
\code{hulk\_vec\_alloc(n,sz)}      & Alloca buffer: header (8 bytes) $+$ $n \times 8$
  bytes de elementos. \\
\code{hulk\_vec\_get(p,i,sz)}      & Devuelve puntero al elemento $i$ con
  verificación de limites. \\
\code{hulk\_vec\_size(p)}          & Lee el header y devuelve el conteo. \\
\code{hulk\_range\_alloc(s,e)}     & Crea iterador para el rango $[s, e)$. \\
\code{hulk\_range\_next(p)}        & Avanza el iterador; devuelve \code{true} si
  hay mas elementos. \\
\code{hulk\_range\_current(p)}     & Valor actual del iterador. \\
\code{hulk\_rand()}                & Numero aleatorio uniforme en $[0, 1)$. \\
\code{hulk\_type\_error(p)}        & Imprime el mensaje y aborta (\code{as} fallido). \\
\bottomrule
\end{tabularx}
\end{table}

\section{Layout en memoria de vectores y rangos}

Las dos estructuras compuestas del \textit{runtime}, \textbf{vectores} y
\textbf{rangos}, tienen layouts compactos diseñados para minimizar el número de
\textit{allocs} y maximizar la localidad de acceso.

\begin{figure}[H]
\centering
\begin{tikzpicture}[font=\small]
  % Vector
  \node at (-0.8,0.4) {\code{Vector}:};
  \draw[fill=hulkblue!8, draw=hulkblue] (0,0) rectangle (2,0.8);
  \draw[fill=halfgray!8, draw=halfgray] (2,0) rectangle (3.5,0.8);
  \draw[fill=halfgray!8, draw=halfgray] (3.5,0) rectangle (5,0.8);
  \draw[fill=halfgray!8, draw=halfgray] (5,0) rectangle (6.5,0.8);
  \draw[fill=halfgray!8, draw=halfgray] (6.5,0) rectangle (8,0.8);
  \node at (1,0.4) {\code{i64 count}};
  \node at (2.75,0.4) {\code{elem[0]}};
  \node at (4.25,0.4) {\code{elem[1]}};
  \node at (5.75,0.4) {\code{elem[2]}};
  \node at (7.25,0.4) {\code{\ldots}};
  \node[font=\scriptsize] at (1,-0.3) {8 bytes};
  \node[font=\scriptsize] at (4.25,-0.3) {8 bytes c.u.};
  % Range
  \node at (-0.8,-1.8) {\code{Range}:};
  \draw[fill=hulkblue!8, draw=hulkblue] (0,-2.2) rectangle (2,-1.4);
  \draw[fill=hulkblue!8, draw=hulkblue] (2,-2.2) rectangle (4,-1.4);
  \draw[fill=hulkaccent!8, draw=hulkaccent] (4,-2.2) rectangle (6,-1.4);
  \node at (1,-1.8) {\code{f64 start}};
  \node at (3,-1.8) {\code{f64 end}};
  \node at (5,-1.8) {\code{f64 current}};
  \node[font=\scriptsize] at (3,-2.5) {24 bytes totales};
\end{tikzpicture}
\caption{Layouts en memoria de vectores y rangos.}
\label{fig:layouts-rt}
\end{figure}

El layout de vector tiene exactamente \textbf{una palabra de cabecera}
($8$ bytes con el conteo) seguida del area de datos. Esto permite calcular el
puntero al elemento $i$ con una única operación de aritmética de punteros:
\code{vec + 8 + 8*i}.

\section{Pipeline AOT completo}

Al invocar el compilador con \code{./hulk source.hulk}, el proceso completo es:

\begin{enumerate}
  \item Análisis léxico, sintáctico y semántico del fuente.
  \item Generación de \llvm{}~IR para todas las declaraciones y la expresión
    principal.
  \item Adición de una función \code{main} compatible con C que llama a
    \code{\_\_hulk\_entry()} y retorna 0.
  \item Optimización del IR con el pipeline descrito en la \cref{tab:passes}.
  \item Emisión de código objeto a un fichero temporal mediante el método
    \code{write\_to\_file} de la \code{TargetMachine} de \llvm{}.
  \item Enlace con \code{gcc}:
    \begin{lstlisting}[style=base,numbers=none,xleftmargin=0pt]
gcc /tmp/hulk_program.o hulk_runtime.a -o ./output -lm
    \end{lstlisting}
\end{enumerate}

El ejecutable resultante es un binario \textbf{ELF x86-64 autocontenido} que no
requiere a \llvm{} en tiempo de ejecución. Sus únicas dependencias son la libc del
sistema y la \code{libm}.

% ════════════════════════════════════════════════════════════════════════════
\chapter{Extensión del lenguaje: vectores con expresiones generadoras}
\label{cap:ext}
% ════════════════════════════════════════════════════════════════════════════

\lettrine[lines=2,findent=2pt,nindent=0pt]{E}{ste} capítulo presenta la extensión
central del lenguaje implementada en este compilador: un \textbf{sistema de
vectores con expresiones generadoras}, estáticamente tipado, integrado con el
sistema de protocolos del lenguaje y con verificación segura de índices en
tiempo de ejecución. La extensión modifica simultáneamente la sintaxis ---añade
literales de vector, una notación generadora inspirada en la \textit{set-builder
notation} matemática, y un operador de indexación---, el sistema de tipos
---introduce un constructor de tipos paramétrico $\mathtt{Vector\langle T\rangle}$
con regla de varianza definida---, la semántica operacional ---introduce reglas
de evaluación con un mecanismo de aborto controlado para índices inválidos--- y
el \textit{runtime} del lenguaje ---añade primitivas de asignación, acceso y
medida de tamaño con verificación de límites.

\section{Motivación}

En su núcleo básico, \hulk{} no incluye estructuras de datos compuestas: las
únicas formas no escalares disponibles son los rangos producidos por la función
\code{range(s,e)} (que tienen una semántica puramente iterable, sin acceso
aleatorio) y los tipos definidos por el usuario. Para escribir algoritmos
elementales que operen sobre colecciones ---calcular el promedio de una lista de
números, aplicar una función a cada elemento, filtrar los valores que cumplen
una condición--- el programador debe declarar un tipo propio que encapsule la
lógica de tamaño, acceso e iteración. Esta declaración no es difícil pero es
sistemáticamente verbosa, en un lenguaje cuya virtud central es precisamente la
concisión expresiva.

La propia especificación del lenguaje reconoce esta carencia y bosqueja, en
una sección dedicada de su apéndice gramatical, un sistema de vectores como
extensión natural. La extensión presentada en este capítulo toma esa propuesta
como punto de partida y la enriquece en cuatro sentidos:

\begin{itemize}
  \item agrega \textbf{indexación} con la sintaxis \code{v[i]} y verificación de
    limites en \textit{runtime};
  \item integra los vectores con el \textbf{protocolo \code{Iterable}}, permitiendo
    usarlos en \code{for} igual que los rangos;
  \item define el \textbf{tipo \code{Vector<T>}} en el sistema de tipos con
    inferencia bidireccional;
  \item proporciona un \textbf{runtime con bounds-checking} que aborta de forma
    controlada en accesos invalidos.
\end{itemize}

\section{Definición formal de la sintaxis}

La extensión agrega tres construcciones sintácticas nuevas a la gramática de
\hulk{}: el vector explicito, la expresión generadora y la indexación.

\begin{lstlisting}[style=bnf, caption={Nuevas producciones gramaticales para la
  extensión.}, label={lst:gram-vec}]
(* Atom es extendido con vectores literales *)
Atom        -> ... | VectorLiteral

VectorLiteral -> '[' ']'                                     (* vector vacio *)
              | '[' ArgListNonEmpty ']'                       (* explicito *)
              | '[' Expr '|' IDENTIFIER 'in' Expr ']'         (* generador *)

(* Postfix es extendido con indexación *)
Postfix     -> Atom (... | '[' Expr ']')*                     (* indexación *)
\end{lstlisting}

\subsection{Vectores explicitos}

Un \textbf{vector explicito} se construye listando sus elementos entre corchetes,
separados por comas. El tipo del vector se infiere a partir del tipo de sus
elementos.

\begin{lstlisting}[style=hulkst,caption={Vectores explicitos.}]
let v     = [1, 2, 3, 4, 5]               in print(v[2]);   // 3
let names = ["Alice", "Bob", "Carol"]     in print(names[0]);
let empty = []                            in print(size(empty)); // 0
\end{lstlisting}

\subsection{Expresiones generadoras}

La sintaxis central de la extensión es la \textbf{expresión generadora}, que crea
un nuevo vector aplicando una expresión a cada elemento de un iterable.

\begin{lstlisting}[style=hulkst,caption={Expresiones generadoras.}]
// Cuadrados de 0 a 4
let squares = [x * x | x in range(0, 5)] in
print(squares[3]);                                 // 9

// Transformación de otro vector
let v       = [1, 2, 3]                  in
let doubled = [x * 2 | x in v]           in
print(doubled[1]);                                 // 4

// El cuerpo puede usar variables del entorno
let factor = 10 in
let scaled = [x * factor | x in range(1, 4)] in
print(scaled[2]);                                  // 30
\end{lstlisting}

La sintaxis es \code{[body | var in iterable]}, donde:
\begin{itemize}
  \item \code{body} es cualquier expresión \hulk{} que puede usar \code{var} y
    variables del entorno circundante;
  \item \code{var} es la variable de iteración, local al cuerpo;
  \item \code{iterable} es un \code{Range} o un \code{Vector<T>} de cualquier
    tipo;
  \item el separador \code{|} coincide con el del operador OR lógico; la
    ambigüedad se resuelve por contexto (vease el conflicto declarado en la
    gramática).
\end{itemize}

\subsection{Indexación}

Los elementos de un vector se acceden con la sintaxis \code{v[i]}, donde \code{i}
es una expresión de tipo \code{Number}.

\begin{lstlisting}[style=hulkst,caption={Indexación con expresión calculada.}]
let v = [10, 20, 30]                                  in
let i = 1                                              in
v[i + 1]   // equivalente a v[2] = 30
\end{lstlisting}

El indice se convierte de \code{Number} (\textit{double}) a entero en tiempo de
compilación mediante \code{fptosi}. El \textit{runtime} verifica que el indice sea
no-negativo y menor que el tamaño del vector; en caso contrario, aborta con un
mensaje descriptivo.

\section{Definición formal de la semántica}

Para precisar el significado de las nuevas construcciones, se presentan las
reglas de tipado y la semántica operacional en estilo natural.

\subsection{Reglas de tipado}

\begin{definition}[title={Definición: tipos de la extensión}]
Se introduce un nuevo constructor de tipos \code{Vector<T>} paramétrico en un
parámetro \code{T}. La conformidad de vectores es \textbf{invariante} en el
parámetro: $\code{Vector<A>} \leq \code{Vector<B>} \iff A = B$.
\end{definition}

Las reglas de tipado para las nuevas expresiones son las siguientes, donde
$\Gamma$ es el ambiente de tipos y $\vdash$ denota el juicio de tipado:

\begin{align*}
  & \frac{\Gamma \vdash e_1 : T \quad \cdots \quad \Gamma \vdash e_n : T}
         {\Gamma \vdash [e_1, \ldots, e_n] : \code{Vector<T>}}
  & \text{(vector explicito)} \\[8pt]
  & \frac{\Gamma \vdash e : \code{Range} \quad \Gamma, x : \code{Number} \vdash b : T}
         {\Gamma \vdash [b \mid x \text{ in } e] : \code{Vector<T>}}
  & \text{(generador sobre rango)} \\[8pt]
  & \frac{\Gamma \vdash e : \code{Vector<S>} \quad \Gamma, x : S \vdash b : T}
         {\Gamma \vdash [b \mid x \text{ in } e] : \code{Vector<T>}}
  & \text{(generador sobre vector)} \\[8pt]
  & \frac{\Gamma \vdash e : \code{Vector<T>} \quad \Gamma \vdash i : \code{Number}}
         {\Gamma \vdash e[i] : T}
  & \text{(indexación)}
\end{align*}

\subsection{Semántica operacional de gran paso}

La semántica de evaluación de los vectores y de los generadores adopta una
estrategia \textbf{ansiosa} (\textit{eager}): todos los elementos se calculan en
orden de izquierda a derecha y se almacenan en un buffer recién asignado
\emph{antes} de que la expresión completa retorne. Esta elección está alineada
con la semántica del resto de \hulk{} ---donde los argumentos de funciones, los
lados de los operadores binarios y los inicializadores de \code{let} se evalúan de
forma estricta--- y se justifica con detalle en la \cref{sec:decisiones-ext}.

Se presenta la semántica en estilo de \textit{gran paso} (también llamada
\textit{natural} o \textit{evaluativa}): el juicio
$\sigma \vdash e \Downarrow v,\,\sigma'$ se lee \textit{``en el entorno $\sigma$,
la expresión $e$ se evalúa al valor $v$ produciendo el entorno $\sigma'$''}. El
entorno $\sigma$ asocia variables a valores, y $\sigma[x \mapsto v]$ denota la
extensión que añade el binding $x \mapsto v$ ocultando cualquier binding previo.

\begin{small}
\[
  \frac{
    \sigma \vdash e_1 \Downarrow v_1,\,\sigma_1
    \quad \cdots \quad
    \sigma_{n-1} \vdash e_n \Downarrow v_n,\,\sigma_n
  }{
    \sigma \vdash [e_1, \ldots, e_n] \Downarrow \mathsf{vec}(v_1, \ldots, v_n),\,\sigma_n
  }
  \tag{E-VecExp}
\]

\[
  \frac{
    \sigma \vdash e \Downarrow \mathsf{vec}(v_1, \ldots, v_n),\,\sigma_0
    \quad
    \forall i.\ \sigma_0[x \mapsto v_i] \vdash b \Downarrow w_i,\,\sigma_i
  }{
    \sigma \vdash [b \mid x \,\text{in}\, e] \Downarrow \mathsf{vec}(w_1, \ldots, w_n),\,\sigma_n
  }
  \tag{E-VecGenV}
\]

\[
  \frac{
    \sigma \vdash e \Downarrow \mathsf{range}(s, t),\,\sigma_0 \quad
    n = \max(0,\, \lfloor t - s \rfloor) \quad
    \forall i \in [0,n).\ \sigma_0[x \mapsto s + i] \vdash b \Downarrow w_i
  }{
    \sigma \vdash [b \mid x \,\text{in}\, e] \Downarrow \mathsf{vec}(w_0, \ldots, w_{n-1}),\,\sigma_n
  }
  \tag{E-VecGenR}
\]

\[
  \frac{
    \sigma \vdash e \Downarrow \mathsf{vec}(v_0, \ldots, v_{n-1}),\,\sigma_0 \quad
    \sigma_0 \vdash i \Downarrow k,\,\sigma_1 \quad 0 \leq \lfloor k \rfloor < n
  }{
    \sigma \vdash e[i] \Downarrow v_{\lfloor k \rfloor},\,\sigma_1
  }
  \tag{E-IndexOk}
\]

\[
  \frac{
    \sigma \vdash e \Downarrow \mathsf{vec}(v_0, \ldots, v_{n-1}),\,\sigma_0 \quad
    \sigma_0 \vdash i \Downarrow k,\,\sigma_1 \quad \lnot(0 \leq \lfloor k \rfloor < n)
  }{
    \sigma \vdash e[i] \Downarrow \bot
  }
  \tag{E-IndexErr}
\]
\end{small}

La regla \textsc{E-IndexErr} captura el comportamiento de aborto controlado del
\textit{runtime}: cuando el índice está fuera de rango, no se produce un valor sino
una transición a un estado de error (notado $\bot$) que el \textit{runtime}
materializa imprimiendo un mensaje en \code{stderr} y llamando a \code{abort(3)}.
Este modelo de fallos sigue el principio de Hoare según el cual \textit{una
condición que no debería ocurrir no debe producir un valor arbitrario; debe
detener la ejecución de la manera más visible posible}~\cite{hoare1985}.

\section{Realización en el código}

\subsection{Modificaciones al lexer}

Ninguna. Los tokens \code{[}, \code{]}, \code{|}, e \code{in} ya existen en el
lexer para otras construcciones. La extensión reutiliza el conjunto existente sin
agregar terminales.

\subsection{Modificaciones al parser}

Se agregan las producciones del \cref{lst:gram-vec}. El único conflicto
shift/reduce nuevo aparece en \code{[Expr | }$\ldots$ porque el token \code{|}
podria iniciar también el operador OR lógico de la regla \code{OrExpr}. El
conflicto se resuelve con un \textit{lookahead} de dos tokens: si tras \code{|}
viene un \code{IDENTIFIER} seguido de \code{in}, se reduce hacia
\code{VectorLiteral}; en cualquier otro caso, se desplaza hacia \code{OrExpr}.

El AST se extiende con dos variantes:

\begin{lstlisting}[style=rustst, caption={Nuevas variantes del AST para la
  extensión.}, label={lst:ast-vec}]
pub enum VectorExpr {
    Explicit  { elements: Vec<Expr>, span: Span },
    Generator { body: Box<Expr>, var: String,
                iterable: Box<Expr>, span: Span },
}

pub struct IndexExpr {
    pub collection: Box<Expr>,
    pub index:      Box<Expr>,
    pub span:       Span,
}
\end{lstlisting}

\subsection{Modificaciones al semántico}

El verificador de tipos implementa las reglas de la sección anterior: para un
vector explicito, infiere el tipo del primer elemento y verifica que los demas
sean compatibles; para una expresión generadora, evalua el tipo del iterable y
construye un ambiente extendido $\Gamma, x : T$ antes de tipar el cuerpo; para
una indexación, verifica que la colección sea \code{Vector<T>} para algun \code{T}
y que el indice sea \code{Number}.

\subsection{Modificaciones al codegen}

La generación de código para vectores explicitos es directa: una llamada a
\code{hulk\_vec\_alloc} para reservar el buffer, seguida de una secuencia de
\code{hulk\_vec\_get} (para obtener punteros a los slots) y \code{store} (para
escribir los valores).

La generación para una expresión generadora es mas elaborada y se ilustra en el
\cref{fig:cfg-gen}: requiere construir un bucle con cuatro bloques básicos
(\code{gen\_cond}, \code{gen\_body}, \code{gen\_end}, mas el bloque de entrada).

\begin{figure}[H]
\centering
\begin{tikzpicture}[
  every node/.style={font=\small},
  bb/.style={rectangle, draw=hulkblue, fill=hulkblue!8, rounded corners=2pt,
    minimum width=4cm, minimum height=0.9cm, align=center},
  arr/.style={-Stealth, thick, color=hulkblue}
]
  \node[bb] (e)  {\textbf{entry}\\\scriptsize compute count, alloc result, init idx=0};
  \node[bb] (c)  [below=of e] {\textbf{gen\_cond}\\\scriptsize idx < count?};
  \node[bb] (b)  [below=of c, xshift=-2.5cm] {\textbf{gen\_body}\\\scriptsize
    fetch elem, eval body, store, idx++};
  \node[bb] (x)  [below=of c, xshift=2.5cm]  {\textbf{gen\_end}\\\scriptsize
    return result};
  \draw[arr] (e) -- (c);
  \draw[arr] (c) -- node[left]{si} (b);
  \draw[arr] (c) -- node[right]{no} (x);
  \draw[arr] (b.east) .. controls (3,-3) and (2,-1.5) .. (c.east);
\end{tikzpicture}
\caption{Grafo de flujo de control del código generado para una expresión
  generadora.}
\label{fig:cfg-gen}
\end{figure}

Una sutileza importante: el calculo del \code{count} para un generador sobre un
\code{Range} hace \code{count = floor(end - start)}. Sin precauciones, un rango con
\code{start > end} produciria \code{count < 0}, lo que provocaria que la
instrucción \code{fptosi} de \llvm{} generara un \textit{poison value} y el
\textit{behavior undefined}. La extensión inserta un \textit{clamp} a
$[0, \textsf{i32::MAX}]$:

\begin{lstlisting}[style=llvmst, caption={Clamp del count para evitar poison en
  fptosi.}, label={lst:clamp}]
%diff   = fsub double %end, %start
%is_neg = fcmp ult double %diff, 0.0
%nn     = select i1 %is_neg, double 0.0, double %diff
%hi     = fcmp ogt double %nn, 2147483647.0
%clamp  = select i1 %hi, double 2147483647.0, double %nn
%count  = fptosi double %clamp to i32
\end{lstlisting}

\section{Análisis comparativo con otros lenguajes}

La construcción implementada esta emparentada con una familia de construcciones
presentes en otros lenguajes bajo nombres como \textit{list comprehension},
\textit{for comprehension} o \textit{generator expression}. A continuación se
compara la extensión propuesta con las equivalentes mas relevantes.

\subsection{Notación matemática de conjuntos}

La inspiración histórica de todas estas construcciones es la \textbf{notación de
construcción de conjuntos} (\textit{set-builder notation}) introducida en la
teoría axiomática de conjuntos a principios del siglo~XX:
$$ S = \{ f(x) \mid x \in A \land P(x) \}. $$

Esta notación, leída \textit{``el conjunto de los $f(x)$ tales que $x$ pertenece a
$A$ y $P(x)$ se cumple''}, separa con una barra vertical la expresión generadora
de los cuantificadores. La extensión propuesta toma directamente esta convención
sintáctica: \code{[body | var in iterable]} es la traducción literal de la
notación matemática al ASCII, con la barra vertical \code{|} como separador y el
\textit{keyword} \code{in} en sustitución del símbolo $\in$. Es importante notar
que la notación matemática puede tener \emph{cualquier} número de generadores y de
predicados; la extensión que aquí se presenta admite, por simplicidad de la
gramática LALR(1), exactamente un generador y ningún predicado. La omisión es
deliberada (vease \cref{sec:decisiones-ext}) y la pérdida de expresividad es
moderada porque cualquier filtro puede emularse insertando un \code{if-else} en el
cuerpo de la comprensión.

\subsection{Comparación con Python}

Python distingue entre \textbf{list comprehensions} (\textit{eager}) y
\textbf{generator expressions} (\textit{lazy}):

\begin{lstlisting}[language=Python,style=base,caption={Comprehensions y generators
  en Python.}]
squares = [x**2 for x in range(5)]      # list comp, eager
gen     = (x**2 for x in range(5))      # generator expr, lazy
\end{lstlisting}

La sintaxis de \hulk{} \code{[body | var in iterable]} corresponde a las
\textit{list comprehensions} de Python tanto en comportamiento (eager,
materializan toda la lista) como en semántica de alcance (\code{var} es local al
cuerpo, no captura el contexto).

Diferencias relevantes:

\begin{itemize}
  \item Python admite \textbf{filtros} (\code{if x > 0}) y \textbf{multiples
    iteradores} (\code{for x in xs for y in ys}); la extensión de \hulk{} se
    limita a un solo iterador sin filtro. Esta simplificación fue intencional:
    mantiene la gramática LALR(1) libre de ambigüedades difíciles y reduce la
    complejidad del codegen.
  \item Python es dinámicamente tipado y permite vectores heterogeneos; \hulk{}
    es estáticamente tipado e impone homogeneidad: todos los elementos deben
    conformar al mismo tipo.
\end{itemize}

\subsection{Comparación con Haskell}

Haskell ofrece \textit{list comprehensions} con una sintaxis aun mas cercana a la
notación matemática:

\begin{lstlisting}[language=Haskell,style=base,caption={List comprehensions en
  Haskell.}]
squares = [x^2 | x <- [0..4]]
evens   = [x | x <- [1..100], even x]
\end{lstlisting}

La evaluación en Haskell es \textbf{perezosa por defecto}: los elementos se
computan solo cuando son demandados. Esto permite construir listas infinitas como
\code{nats = [0..]}, donde los valores se calculan incrementalmente.

\hulk{} adopta evaluación \textbf{ansiosa}. Las razones de esta decisión se
exponen en la \cref{sec:decisiones-ext}. El precio que se paga es la
imposibilidad de manejar colecciones infinitas; el beneficio es un modelo de
memoria y de ejecución mucho mas simple, compatible con el resto del lenguaje
(que también es eager).

\subsection{Comparación con Scala}

Scala ofrece la construcción \code{for...yield}, una azucar sintáctica sobre
\code{flatMap} y \code{map}:

\begin{lstlisting}[language=Scala,style=base,caption={For-comprehension en Scala.}]
val squares = for (x <- 0 until 5) yield x * x
val pairs   = for {
  x <- 1 to 3
  y <- 1 to x
} yield (x, y)
\end{lstlisting}

Scala es \textit{eager} para colecciones concretas como \code{Vector} y
\textit{lazy} para vistas \code{.view}. La sintaxis es mas verbosa pero soporta
multiples generadores y guardas. La elección de \hulk{} de exigir un solo
generador minimiza la complejidad sintáctica al precio de la expresividad.

\subsection{Comparación con C++20 ranges}

C++20 introdujo el sistema de \textbf{ranges y views} con evaluación perezosa y
composición por \textit{pipe}:

\begin{lstlisting}[language=C++,style=base,caption={Ranges y views en C++20.}]
auto squares = std::views::iota(0, 5)
             | std::views::transform([](int x){ return x*x; });
std::vector<int> v(squares.begin(), squares.end());  // materializa
\end{lstlisting}

C++ favorece evaluación perezosa con \textit{materialización explicita}, lo que
permite optimizar la composición de transformaciones evitando allocaciones
intermedias. Para \hulk{} --- que no posee \textit{iterators} genéricos ni un
sistema de tipos con \textit{lifetimes} --- esta aproximación seria
desproporcionadamente compleja.

\subsection{Resumen comparativo}

\begin{table}[H]
\centering
\caption{Comparativa de comprehensions/generators entre lenguajes.}
\label{tab:comp-lang}
\begin{tabularx}{\textwidth}{lXcXc}
\toprule
\textsc{Lenguaje} & \textsc{Sintaxis} & \textsc{Evaluación}
  & \textsc{Filtros} & \textsc{Tipos} \\
\midrule
HULK    & \code{[e | x in it]}            & Ansiosa      & No  & Estático inferido \\
Python  & \code{[e for x in it if c]}     & Ansiosa/Diferida   & Sí  & Dinámico \\
Haskell & \code{[e | x <- xs, c]}         & Diferida     & Sí  & Estático inferido \\
Scala   & \code{for x <- xs yield e}      & Ansiosa/Diferida   & Sí  & Estático inferido \\
C++20   & \code{xs | transform(f)}        & Diferida     & Sí  & Estático explícito \\
\bottomrule
\end{tabularx}
\end{table}

\section{Decisiones de diseño de la extensión}
\label{sec:decisiones-ext}

\subsection{Por que evaluación ansiosa y no perezosa}

Tres argumentos sostienen la decisión de evaluación ansiosa.

\paragraph{Simplicidad del modelo de memoria.}
Con evaluación ansiosa el tamaño del vector se conoce antes de la asignación,
por lo que se requiere una sola llamada a \code{malloc} con la capacidad exacta.
La evaluación perezosa requeriria \textit{thunks} (cierres con estado) o
continuaciones, estructuras adicionales que duplicarian la complejidad del
\textit{runtime} sin un beneficio significativo en el caso de uso típico.

\paragraph{Rendimiento predecible.}
Los compiladores de lenguajes perezosos como GHC para Haskell invierten esfuerzo
considerable en evitar los \textit{overheads} de la \textit{thunkificación}: el
análisis de demanda, el análisis de \textit{strictness} y la deforestación son
optimizaciones complejas dedicadas a este problema. En un compilador con fines
educativos que ya delega la optimización en \llvm{}, la evaluación ansiosa
produce IR simple y fácil de optimizar con los \textit{passes} estandar.

\paragraph{Consistencia con el resto del lenguaje.}
\hulk{} tiene efectos de lado (\code{print}, asignación destructiva). La
evaluación perezosa en presencia de efectos produce comportamientos sorprendentes
porque el orden de evaluación deja de coincidir con el orden textual. La
evaluación ansiosa garantiza que los efectos se producen de izquierda a derecha,
en consonancia con la semántica del resto de \hulk{}.

\subsection{Por que la barra simple \code{|} como separador}

La barra vertical reproduce la notación matemática $\{f(x) \mid x \in S\}$ y
coincide con la convención de Haskell \code{[f x | x <- xs]}. El hecho de que
\code{|} sea también el operador OR lógico introduce un conflicto shift/reduce
que se resuelve con un \textit{lookahead} adicional. Alternativas consideradas:

\begin{itemize}
  \item \code{[body for var in iterable]} --- mas familiar para programadores
    Python, pero introduce el keyword \code{for} en un contexto distinto al de
    los bucles \code{for}, lo que dificulta el aprendizaje del lenguaje.
  \item \code{[body : var in iterable]} --- colisiona con la notación de tipo
    \code{x : T} ya presente en la gramática de \hulk{}.
  \item \code{[body || var in iterable]} --- evita el conflicto pero rompe la
    correspondencia con la notación matemática.
\end{itemize}

\subsection{Por que vectores inmutables (tamaño fijo)}

Los vectores de \hulk{} no tienen operaciones de \code{push} o \code{pop}: una vez
creados, su tamaño no cambia. Las modificaciones afectan solo a los elementos.
Esta decisión se justifica por:

\begin{enumerate}
  \item \textbf{Simplicidad del runtime}: un buffer fijo con cabecera es trivial.
    Vectores dinámicos requeririan estrategias de redimensión (duplicar
    capacidad) que duplicarian la complejidad de la biblioteca de soporte.
  \item \textbf{Semántica de expresión}: en \hulk{} todo es una expresión. Un
    vector dinámico con mutación es un objeto con estado, lo que rompe la
    composición funcional que permite escribir
    \code{print([x*x | x in range(0,5)][3])}.
  \item \textbf{Caso de uso del lenguaje}: \hulk{} es un lenguaje educativo;
    sus programas típicos son cortos y operan sobre vectores de tamaño conocido.
\end{enumerate}

% ════════════════════════════════════════════════════════════════════════════
\chapter{Cobertura de la especificación}
\label{cap:cov}
% ════════════════════════════════════════════════════════════════════════════

\lettrine[lines=2,findent=2pt,nindent=0pt]{E}{ste} capítulo cuantifica el grado de
cobertura del compilador respecto a la especificación oficial del lenguaje
\hulk{}. La especificación organiza el lenguaje en las secciones A.2 a A.14 de
su apéndice gramatical: las secciones A.2 a A.8 forman el núcleo del lenguaje,
mientras que las restantes describen extensiones opcionales que pueden
incorporarse de manera incremental.

\section{Tabla de cobertura por sección}

\begin{table}[H]
\centering
\caption{Cobertura de la especificación oficial de \hulk{} por sección.}
\label{tab:cobertura}
\begin{tabularx}{\textwidth}{clcc}
\toprule
\textsc{Sec.} & \textsc{Característica} & \textsc{Obl.} & \textsc{Estado} \\
\midrule
A.2.1 & Expresiones aritméticas, precedencia & si & \done \\
A.2.2 & Cadenas, concatenación, escapes      & si & \done \\
A.2.3 & Funciones matemáticas builtin        & si & \done \\
A.2.4 & Bloques de expresiones               & si & \done \\
A.3   & Funciones (inline y completas)       & si & \done \\
A.4   & Variables y \code{let}-\code{in}     & si & \done \\
A.4.6 & Asignación destructiva               & si & \done \\
A.5   & Condicionales \code{if}-\code{elif}-\code{else} & si & \done \\
A.6.1 & Ciclo \code{while}                   & si & \done \\
A.6.2 & Ciclo \code{for}                     & si & \done \\
A.7   & Tipos, herencia, \code{self}, \code{base()} & si & \done \\
A.8.1 & Anotaciones de tipo                  & si & \done \\
A.8.2 & Conformidad \code{<=}                & si & \done \\
A.8.3 & Operador \code{is}                   & si & \done \\
A.8.4 & Operador \code{as} (downcast)        & si & \done \\
\midrule
A.9   & Inferencia de tipos (parcial)        & no & \halfdone \\
A.10  & Protocolos y conformidad estructural & no & \done \\
A.11  & Iterables (\code{Iterable}/\code{Range})  & no & \done \\
A.12  & \textbf{Vectores y generadores (extensión)} & no & \done \\
A.13  & Funtores y lambdas                   & no & \miss \\
A.14  & Macros con pattern matching          & no & \miss \\
\midrule
\multicolumn{2}{r}{\textsc{Cobertura obligatoria}}  & \multicolumn{2}{c}{\textbf{100\%}} \\
\bottomrule
\end{tabularx}
\end{table}

\section{Ejemplos representativos}

El compilador es capaz de procesar programas como los siguientes, que ejercitan
multiples características a la vez.

\subsection{Herencia y polimorfismo}

\begin{lstlisting}[style=hulkst, caption={Herencia, polimorfismo y vectores
  heterogeneos.}]
type Animal(name: String) {
    name = name;
    speak(): String => "...";
    greet(): String =>
        "I am " @@ self.name @@ " and I say: " @@ self.speak();
}

type Dog(name: String) inherits Animal(name) {
    speak(): String => "Woof!";
}

type Cat(name: String) inherits Animal(name) {
    speak(): String => "Meow!";
}

let zoo = [new Dog("Rex"), new Cat("Whiskers")] in
for (a in zoo) print(a.greet());
\end{lstlisting}

\subsection{Protocolos con conformidad estructural}

\begin{lstlisting}[style=hulkst, caption={Protocolos y conformidad implicita.}]
protocol Printable {
    to_str(): String;
}

type Point(x: Number, y: Number) {
    x = x; y = y;
    to_str(): String => "(" @ x @ ", " @ y @ ")";
}

function show(item: Printable) => print(item.to_str());
show(new Point(3, 4));
\end{lstlisting}

\subsection{Vectores y generadores}

\begin{lstlisting}[style=hulkst, caption={Calculo de los primeros 10 Fibonacci con
  generadores.}]
function fib(n: Number): [Number] =>
    let result = [0 | i in range(0, n)] in {
        result[0] := 0; result[1] := 1;
        for (i in range(2, n)) result[i] := result[i-1] + result[i-2];
        result
    };

for (x in fib(10)) print(x);
\end{lstlisting}

% ════════════════════════════════════════════════════════════════════════════
\chapter{Decisiones de diseño significativas}
\label{cap:dec}
% ════════════════════════════════════════════════════════════════════════════

\lettrine[lines=2,findent=2pt,nindent=0pt]{T}{odo} compilador real es una secuencia
de decisiones de compromiso entre simplicidad de implementación, eficiencia del
código generado y completitud del lenguaje soportado. Este capítulo documenta
las decisiones más significativas tomadas durante la construcción del compilador
y articula explícitamente las alternativas descartadas, los argumentos a favor de
la opción elegida y los costos asumidos.

\section{LALR(1) artesanal frente a generador externo}

El parser está implementado desde cero ---gramática como estructura de datos,
tablas LALR(1) construidas programáticamente en tiempo de inicialización--- en
lugar de generarse a partir de YACC, Bison, LALRPOP o ANTLR.

\begin{description}[style=nextline]
  \item[Argumento principal]
    \textit{Transparencia algorítmica.} La construcción del autómata LR(0), el
    cómputo de FIRST y FOLLOW, la propagación de \textit{lookaheads} y la
    construcción de las tablas \textsc{Action}/\textsc{Goto} constituyen el núcleo
    técnico de la teoría de \textit{parsing}~\cite{aho2006,deremer1969}. Una
    implementación que delegue ese núcleo en una herramienta externa convierte
    al \textit{parser} en una caja negra, lo que va en contra del principio de
    diseño D1 enunciado en la introducción.
  \item[Argumento secundario]
    \textit{Inspeccionabilidad y depurabilidad.} La gramática es una estructura
    de datos en memoria que puede imprimirse, recorrerse y analizarse en
    tiempo de ejecución. Esto resulta determinante a la hora de diagnosticar
    conflictos \textit{shift}/\textit{reduce}: el conflicto único de la
    gramática (asociado al separador \code{|} de las expresiones generadoras)
    pudo identificarse y resolverse inspeccionando la tabla generada, sin
    necesidad de recompilar el compilador entre iteraciones.
  \item[Coste]
    Implementación mas larga (aproximadamente 800 LOC adicionales solo para el
    modulo LALR) y mayor tiempo de \textit{startup} debido a que las tablas se
    construyen en cada invocación. Este \textit{overhead} es de pocos
    milisegundos en programas típicos.
\end{description}

\section{LLVM como backend vs interprete vs VM propia}

La especificación permite generar código máquina directamente, usar LLVM, o
implementar una máquina virtual propia.

\begin{description}[style=nextline]
  \item[Argumento principal]
    \textit{Calidad del código generado.} \llvm{} produce código comparable al
    de \code{gcc -O2} para estructuras comunes. Una máquina virtual propia o un
    interprete serian varios ordenes de magnitud mas lentos para programas con
    aritmética intensiva.
  \item[Argumento secundario]
    \textit{Portabilidad.} El mismo IR de \llvm{} puede compilarse para x86-64,
    ARM, RISC-V o WebAssembly cambiando un solo parámetro en
    \code{TargetMachine}.
  \item[Coste]
    Curva de aprendizaje de \llvm{} y dependencia de \code{libLLVM-17.so} en el
    sistema de \textit{build}. El paquete \inkwell{} mitiga la primera dificultad.
\end{description}

\section{Vtables vs tagged union vs boxed types}

Para implementar el polimorfismo dinámico se evaluaron tres alternativas:

\begin{enumerate}
  \item \textbf{Tagged union.} Cada objeto es una union discriminada con un tag
    de tipo. Simple pero ineficiente: el tamaño de cada objeto es el del tipo
    mayor de la jerarquía.
  \item \textbf{Boxed types con dispatch dinámico por nombre.} Cada llamada
    a método consulta una tabla \textit{hash} por nombre del método. Funciona
    pero introduce \textit{overheads} de \textit{hashing} y posibles colisiones.
  \item \textbf{Vtables} (\textit{elegida}). Cada objeto tiene un puntero a una
    tabla de funciones. Es la técnica usada por C++, Java JIT, y Rust (para
    \textit{trait objects}). El costo por llamada es una carga adicional desde
    cache; el beneficio es el layout estático y conocido.
\end{enumerate}

\section{Type tags DFS pre-orden vs class objects}

El test de tipo \code{is} se podria implementar de varias formas:

\begin{enumerate}
  \item \textbf{Recorrer la cadena de herencia} en \textit{runtime} comparando
    nombres o punteros a \textit{class objects}. $O(h)$ donde $h$ es la altura
    de la jerarquía.
  \item \textbf{Tabla de \textit{class objects}}. Almacenar en cada objeto un
    puntero a su \textit{class}, y consultar una tabla precomputada de
    relaciones. $O(1)$ con espacio cuadrático en el número de tipos.
  \item \textbf{Type tags DFS pre-orden} (\textit{elegida}). $O(1)$ en tiempo,
    $O(N)$ en espacio.
\end{enumerate}

La elección DFS pre-orden es la combinación óptima tiempo-espacio para el caso
común.

\section{Enlace dinámico vs estático de \llvm{}}

El compilador enlaza dinámicamente con \code{libLLVM-17.so}. La alternativa
(enlace estático) produciría ejecutables de 50--200~MB. El enlace dinámico
mantiene el binario en pocos megabytes y comparte la biblioteca con otras
herramientas del sistema.

% ════════════════════════════════════════════════════════════════════════════
\chapter{Evaluación y pruebas}
\label{cap:ev}
% ════════════════════════════════════════════════════════════════════════════

\lettrine[lines=2,findent=2pt,nindent=0pt]{L}{a} validación del compilador se
realiza mediante una bateria de pruebas automatizadas que cubren todas las fases
y todas las construcciones del lenguaje. Esta sección describe la estructura de
las pruebas y los resultados obtenidos.

\section{Estructura de la suite de pruebas}

\begin{table}[H]
\centering
\caption{Distribución de las pruebas automatizadas por modulo.}
\label{tab:pruebas}
\begin{tabularx}{\textwidth}{lrX}
\toprule
\textsc{Modulo} & \textsc{Tests} & \textsc{Cobertura} \\
\midrule
Lexer      & 87  & Todos los tipos de token, edge cases, escapes \\
Parser     & 112 & Todos los no-terminales, precedencia, asociatividad \\
Semántico  & 98  & Inferencia, protocolos, errores de tipo \\
Codegen    & 87  & Aritmética, flujo, OOP, vectores, rangos, builtins \\
\midrule
\textbf{Total} & \textbf{384} & \\
\bottomrule
\end{tabularx}
\end{table}

\section{Programas de prueba completos}

Mas alla de las pruebas unitarias, el repositorio incluye 22 programas
\hulk{} completos en \code{tests/hulk/}, organizados por categoría:

\begin{itemize}
  \item \code{ok/minimal/}: 9 programas (aritmética, condicionales, funciones,
    hola mundo, \code{let}-\code{in}, cadenas, ciclos).
  \item \code{ok/oop/}: 3 programas (clase básica, herencia, mutación).
  \item \code{ok/types/}: 3 programas (anotaciones, builtins, inferencia).
  \item \code{ok/extras/}: 2 programas (for, range).
  \item \code{errors/}: 5 programas que deben fallar con un mensaje específico.
\end{itemize}

Cada programa se compila y ejecuta, comparando la salida estandar contra una
salida esperada. Los programas de error verifican adicionalmente que el código
de salida sea distinto de cero y que el mensaje siga el formato del contrato de
interfaz.

\section{Contrato de errores}

El compilador respeta el contrato de interfaz asociado al lenguaje: los errores
se reportan por \code{stderr} con el prefijo \code{!\ } y el formato
\code{(línea,col) FASE: mensaje}; el proceso termina con código de salida $\neq 0$.

% ════════════════════════════════════════════════════════════════════════════
\chapter{Limitaciones conocidas y trabajo futuro}
\label{cap:lim}
% ════════════════════════════════════════════════════════════════════════════

\lettrine[lines=2,findent=2pt,nindent=0pt]{N}{ingun} compilador es completo. Este
capitulo documenta las limitaciones identificadas en la implementación actual y
las direcciones de trabajo futuro consideradas.

\section{Limitaciones actuales}

\paragraph{Funciones anónimas (lambdas).}
La especificación (sección~A.13) describe lambdas como funciones de primera
clase. La implementación actual solo soporta funciones con nombre. La adición de
lambdas requeriría \textit{closures} con captura del entorno léxico, lo que
implica un cambio en la representación de los valores funcionales: cada lambda
debería representarse como un par
\code{(puntero\ a\ función,\ puntero\ a\ entorno\ capturado)}, con la
correspondiente gestión de tiempo de vida sobre el entorno. La adición es
mecánica pero no trivial: el código actual del codegen y el del semántico están
construidos bajo el supuesto de que toda función es global, y ese supuesto
permea muchas estructuras de datos.

\paragraph{Procesamiento de secuencias de escape.}
Como se anticipó en el análisis léxico, las secuencias \code{\textbackslash n},
\code{\textbackslash t}, \code{\textbackslash"} y \code{\textbackslash\textbackslash}
no se descodifican y aparecen literalmente en la salida. La implementación es
breve (un bucle sobre los caracteres del lexema), pero requiere decidir
cuidadosamente dónde se realiza la descodificación: hacerlo en el \textit{lexer}
desincroniza la posición del lexema con la del fuente original, lo que dañaría la
calidad de los mensajes de error; hacerlo en la fase semántica preserva la
posición pero implica recorrer el lexema dos veces (una para descodificar, otra
para construir el literal). Esta es una de las tareas pendientes de menor
complejidad y mayor visibilidad para el usuario.

\paragraph{Gestión de memoria.}
El compilador adopta una estrategia de \textit{solo-asignación}: las cadenas y
los objetos se asignan en \textit{heap} con \code{malloc} y nunca se liberan.
Para programas con muchas operaciones de concatenación o asignación de objetos,
esto produce un crecimiento monótono del consumo de memoria. Una mejora futura
podría ser un \textit{arena allocator} de \textit{regiones} ligado al
\code{let}-\code{in} ---libera todas las cadenas creadas dentro de un \code{let}
al salir de él---, o un sistema de \textit{reference counting} con
detección de ciclos similar al de Python.

\paragraph{Recuperación de errores en el parser.}
El parser se detiene en el primer error sintáctico. Un parser con recuperación
podría reiniciar en el siguiente punto de sincronización (\code{;} o \code{\}}) y
reportar todos los errores en una sola pasada. La técnica clásica de Burke y
Fisher~\cite{cooper2011} es directamente aplicable a la tabla generada por el
constructor LALR(1) sin necesidad de regenerar la gramática.

\paragraph{Spans en errores de codegen.}
Los errores internos del generador de código no incluyen la posición en el
fuente. Las fases anteriores propagan \code{Span} en todos los nodos del AST;
propagarlos hasta el codegen mejoraría la experiencia de diagnóstico en caso de
errores internos del compilador.

\paragraph{Vectores sin mutación estructural.}
Los vectores son de tamaño fijo: una vez creados, no admiten \code{push} ni
\code{pop}. Agregar mutación estructural requeriría sustituir el layout actual
(un puntero al buffer con cabecera de tamaño) por una estructura de
\textit{capacity-and-length} similar al \code{Vec<T>} de \rust{}, con una
estrategia de redimensión exponencial para amortizar el coste.

\section{Trabajo futuro}

\begin{itemize}
  \item \textbf{Closures y lambdas}: representar funciones de primera clase con
    un par (\textit{puntero a función}, \textit{puntero a entorno}).
  \item \textbf{Generadores perezosos}: variantes lazy de las expresiones
    generadoras para permitir colecciones conceptualmente infinitas.
  \item \textbf{Inferencia Hindley-Milner completa}: el sistema actual infiere
    tipos a traves de expresiones simples pero no resuelve constraints
    polimorficos complejos. H-M permitiria funciones genéricas como
    \code{function map(v: [A], f: A -> B): [B]}.
  \item \textbf{Optimizaciones específicas de HULK}: \textit{inlining} de
    métodos monomorficos, análisis de escape para objetos de corta vida
    (sustituyendo \textit{heap} por \textit{stack}).
  \item \textbf{Backend WebAssembly}: cambiar el \code{TargetMachine} de
    \llvm{} para emitir WebAssembly permitiria ejecutar programas \hulk{} en el
    navegador.
\end{itemize}

% ════════════════════════════════════════════════════════════════════════════
\chapter{Conclusiones}
\label{cap:concl}
% ════════════════════════════════════════════════════════════════════════════

\lettrine[lines=2,findent=2pt,nindent=0pt]{E}{l} compilador descrito en este reporte
implementa el lenguaje \hulk{} de extremo a extremo: a partir de un fichero
fuente \texttt{.hulk}, produce un binario ELF~x86-64 que el sistema operativo
carga y ejecuta sin más dependencias que \code{libc}, \code{libm} y la
biblioteca de \textit{runtime} aportada por la propia distribución. Las
características nucleares del lenguaje ---expresiones aritméticas y de cadena,
control de flujo expresivo, funciones de primer orden, tipos con herencia
simple, protocolos con conformidad estructural variante y operadores
\code{is}/\code{as} de prueba y conversión de tipos--- quedan implementadas en
su totalidad, junto al sistema de iterables y la extensión propia de vectores
con expresiones generadoras desarrollada en el \cref{cap:ext}.

Tres decisiones técnicas estructuran la implementación. La primera ---construir
el \textit{frontend} a mano en lugar de generarlo a partir de una herramienta
externa--- privilegia la transparencia del autómata LALR(1) y permite resolver
en el constructor de la tabla un conflicto deliberado que ninguna herramienta
genérica acomoda con elegancia (el separador \code{|} de las expresiones
generadoras, idéntico al operador OR lógico). La segunda ---asignar etiquetas
de tipo en orden DFS pre-orden sobre la jerarquía de herencia--- reduce el
operador \code{is} a una comparación entera de dos cláusulas, una mejora que
sitúa el coste del despacho de tipos por debajo del de cualquier comprobación
implementada como recorrido de tabla. La tercera ---resolver el despacho a
través de protocolos mediante un \textit{switch} sobre el \code{type\_tag}, con
ramas terminadas en llamadas estáticas--- contrasta con el mecanismo habitual
en C++ y Rust (tablas adicionales por interfaz) y aprovecha el hecho de que el
verificador semántico conoce en tiempo de compilación el conjunto completo de
tipos conformes a cada protocolo invocado.

La extensión de vectores con expresiones generadoras sitúa a \hulk{} en la
familia de lenguajes que ofrecen soporte de primera clase para transformaciones
declarativas de colecciones, junto a Python, Haskell, Scala y, más
recientemente, C++20. Las diferencias respecto a esos cuatro lenguajes ---y
documentadas con detalle en la \cref{tab:comp-lang}--- no son tanto de
expresividad como de \emph{integración}: la extensión convive con el sistema
de tipos estático del lenguaje, con los protocolos estructurales y con la
jerarquía de objetos sin introducir un subsistema paralelo. El precio asumido
ha sido renunciar a filtros, a generadores múltiples y a evaluación perezosa;
el beneficio, un \textit{runtime} de pocos cientos de bytes y un IR \llvm{}
que el \textit{pass manager} optimiza con la misma eficacia que un bucle
aritmético equivalente escrito a mano. Cada uno de esos compromisos está
justificado por separado en la \cref{sec:decisiones-ext}.

\vspace{1em}
Las direcciones de trabajo enumeradas en el \cref{cap:lim} delimitan las
extensiones naturales del compilador en su estado actual: la introducción de
funciones de primera clase con cierres léxicos, el procesamiento de las
secuencias de escape en literales de cadena, la recuperación de errores
sintácticos para reportar múltiples errores en una sola compilación, y la
sustitución del esquema actual de asignación monotónica por un
\textit{allocator} basado en regiones ligadas al ámbito de las expresiones
\code{let}. Ninguna de ellas es conceptualmente difícil; cada una representa
una ampliación localizada del compilador construido.

% ════════════════════════════════════════════════════════════════════════════
\begin{thebibliography}{99}
\addcontentsline{toc}{chapter}{Bibliografía}
% ════════════════════════════════════════════════════════════════════════════

\bibitem{aho2006}
A.~V. Aho, M.~S. Lam, R.~Sethi, J.~D. Ullman.
\textit{Compilers: Principles, Techniques, and Tools} (2nd ed.).
Addison-Wesley, 2006.

\bibitem{thompson1968}
K.~L. Thompson.
``Programming Techniques: Regular Expression Search Algorithm.''
\textit{Communications of the ACM}, 11(6):419--422, 1968.

\bibitem{deremer1969}
F.~L. DeRemer.
\textit{Practical Translators for LR(k) Languages}.
PhD thesis, MIT, 1969.

\bibitem{lattner2004}
C.~Lattner, V.~Adve.
``LLVM: A Compilation Framework for Lifelong Program Analysis and Transformation.''
\textit{Proceedings of the International Symposium on Code Generation and
Optimization}, 2004.

\bibitem{bjarne1994}
B.~Stroustrup.
\textit{The Design and Evolution of C++}.
Addison-Wesley, 1994.

\bibitem{pierce2002}
B.~C. Pierce.
\textit{Types and Programming Languages}.
MIT Press, 2002.

\bibitem{peytonjones2003}
S.~Peyton Jones (ed.).
\textit{Haskell 98 Language and Libraries: The Revised Report}.
Cambridge University Press, 2003.

\bibitem{odersky2008}
M.~Odersky, L.~Spoon, B.~Venners.
\textit{Programming in Scala}.
Artima Press, 2008.

\bibitem{klabnik2019}
S.~Klabnik, C.~Nichols.
\textit{The Rust Programming Language}.
No Starch Press, 2019.

\bibitem{inkwell}
inkwell: Safe Rust bindings for LLVM.
\url{https://github.com/TheDan64/inkwell}.

\bibitem{bringhurst2004}
R.~Bringhurst.
\textit{The Elements of Typographic Style} (3rd ed.).
Hartley \& Marks, 2004.

\bibitem{miede2011}
A.~Miede.
\textit{A Classic Thesis Style: An Homage to The Elements of Typographic Style}.
LaTeX template, 2011.

\bibitem{appel1998}
A.~W. Appel.
\textit{Modern Compiler Implementation in ML}.
Cambridge University Press, 1998.

\bibitem{cooper2011}
K.~D. Cooper, L.~Torczon.
\textit{Engineering a Compiler} (2nd ed.).
Morgan Kaufmann, 2011.

\bibitem{hoare1985}
C.~A.~R. Hoare.
\textit{Communicating Sequential Processes}.
Prentice Hall International, 1985.

\bibitem{nystrom2021}
R.~Nystrom.
\textit{Crafting Interpreters}.
Genever Benning, 2021.
\url{https://craftinginterpreters.com/}

\bibitem{wadler1990}
P.~Wadler.
``Comprehending Monads''.
\textit{Proceedings of the 1990 ACM Conference on LISP and Functional
Programming}, pp.~61--78, 1990.

\bibitem{wirth1971}
N.~Wirth.
``The Design of a Pascal Compiler''.
\textit{Software: Practice and Experience}, 1(4):309--333, 1971.

\end{thebibliography}

\end{document}
